\chaptertitle{Structure and Strangeness}

                                                                         Section IV

                                            Structure and Strangeness
        Mathematical structures are among the most beautiful discoveries made by the
human mind. The best of these discoveries have tremendous metaphorical and
explanatory power, jumping across' discipline boundaries, illuminating many areas of
thought simultaneously. In addition, the best discoveries often reveal truly bizarre
facets of familiar concepts. In the following seven chapters, four wonderful
mathematical ideas are considered. The "Magic Cube" is an engaging object for many
reasons, not the least of which is its seeming physical impossibility, as well as the
frustrating way that order and chaos appear and disappear on its surface as it is
twisted. The borderline between order and chaos in mathematics is the next topic
treated,- where we see the iteration of very simple functions giving rise to
unexpectedly chaotic phenomena-in particular, "strange attractors". A strange
attractor is a very peculiar shape having structure on an infinite number of scales at
once. This property, applies not only to strange attractors, but to a much larger class
of shapes known as "fractals". They in turn are examples of the more general
mathematical concept of recursion, one of our era's most fruitful areas of exploration
in mathematics and computer science. Recursion and recursivity are presented in three
chapters on the computer language Lisp, the language used most in artificial
intelligence research. Finally, we move from computers to their microscopic -
substrate: the eerie netherworld of quantum phenomena, and the unresolved mysteries
about the relationship between the macroworld and the microworld.




Magic Cubology                                                                     300


                                Magic Cubology
                                     March, 1981

                                                Cubitis magikia, n. A severe mental
                                                disorder accompanied by itching of the
                                                fingertips, which can be relieved only
                                                by prolonged contact with a
                                                multicolored cube originating in
                                                Hungary and Japan. Symptoms often
                                                last for months. Highly contagious.


WHAT this stuffy medical-dictionary entry fails to mention is that contact with
the multicolored cube not only cures the itchiness but also causes it. Furthermore, it
fails to point out that the affliction can be highly pleasurable. I ought to know; I have
suffered from it for the past year and still exhibit the symptoms.
         Búvás Kocka--the Magic Cube, also known as Rubik's Cube-has
simultaneously taken the puzzle world, the mathematics world, and the computing
world by storm. (See Figure 14-1.) Seldom has a puzzle so fired the imagination of so
many people, perhaps not since Sam Lloyd’s famous "15" Puzzle, which caused mass
insanity when it came out in the nineteenth century, and which is still one of the
world's most popular puzzles. The 15 Puzzle and the Magic Cube are spiritual kin, the
one being a two-dimensional problem of restoring the scrambled numbered pieces of
a 4X4 square to their proper positions, and the other being a three-dimensional
problem of restoring the scrambled colored pieces of a 3 X 3 X 3 cube to their proper
positions. The solutions of both demand that the solver be willing to undo seemingly
precious progress time and time again; there is no route to the goal that does not call
for partial but temporary destruction of the visible order achieved up to a given point.
If this is a-difficult lesson to learn with the 15 Puzzle, how much harder with the
Magic Cube! And both puzzles have the fiendish property that well-meaning
bumblers or cunning rogues can take them apart and put them back together in
innocent-looking positions from which the goal is




Magic Cubology                                                                       301

FIGURE 14-1. A Magic Cube in (a) its pristine date, also called START; (b) a typical
scrambled state.

absolutely unattainable, thereby causing the would-be solver considerable
consternation.
        This Magic Cube is much more than just a puzzle. It is an ingenious
mechanical invention, a pastime, a learning tool, a source of metaphors, an
inspiration. It now seems an inevitable object, but it took a long time to be discovered.
Somehow, though, the time was ripe, because the ides germinated and developed
nearly in parallel in Hungary and Japan and perhaps even elsewhere. A report
surfaced recently of a French inspector general named Semah, who claims to
remember encountering such a cube made out of wood in 1920 in Istanbul and then
again in 1935 in Marseilles. Of course, without confirmation the claims seem dubious,
but still titillating. In any event, Rubik's work was completed by 1975, and his
Hungarian patent bears that date. Quite independently, Terutoshi Ishige, a self-taught
engineer and the owner of a small ironworks near Tokyo, carne up with much the
same design within a year of Rubik and filed for a Japanese patent in 1976. Ishige
also deserves credit for this wonderful insight.
        Who is Rubik? Ernó Rubik is a teacher of architecture and design at the
School for Commercial Artists in Budapest. Seeking to sharpen his students' ability to
visualize three-dimensional objects, he came up with the idea of a 3 X 3 X 3 cube any
of whose six 3 X 3 faces could rotate about its center, yet in such a way that the cube
as a whole would not fall apart. Each face would initially be colored uniformly, but
repeated rotations of the various faces would scramble the colors horribly. Then his
students had to figure out how to undo the scrambling.
        When I first heard the cube described over the telephone, it sounded like a
physical impossibility. By all logic, it ought to fall apart into its constituent "cubies"
(one of the many useful and amusing terms invented by "cubists" around the world).
Take any corner cubie-what is it attached to? By imagining rotating each of the three
faces to which it belongs, you can see that the corner cubie in question is detachable
from each of its three




Magic Cubology                                                                        302

edge-cubie neighbors. So how in the world is it held in place? Some people postulate
magnets, rubber bands, or elaborate systems of twisting wires in the interior of the
cube, yet the design is remarkably simple and involves no such items.
        In fact, the Magic Cube can be disassembled in a few seconds (see Figure 14-
2c), revealing an infernal structure so simple that one has to ponder how it can do
what it does. To see what holds it together, first observe that there are three types of
cubie: six center cubies, twelve edge cubies, and eight corner cubies. (See Figure 14-
2a.) Each center cubie has only one "facelet"; edge cubies have two; corner cubies
have three. Moreover, the six center cubies are really not cubical at all-they are just
square façades covering the tips of axles that sprout out from a sixfold spindle in the
cube's heart. The other cubies, however, are nearly complete little cubes, except that
each one has a blunt little "foot" reaching toward the middle of the cube, and some
curved nicks facing inward.
        The basic trick is that cubies mutually hold one another in by means of their
feet, without any cubie actually being attached to any other. Edge cubies hold corner
cubies' feet, corner cubies hold edge cubies' feet. Center cubies are the keystones. As
any layer, say the top one, rotates, it holds itself together horizontally, and is held in
place vertically by its own center and by the equatorial layer below it. The equatorial
layer has a sunken circular track (formed by the nicks in its cubies) that guides the
motion of the upper layer's feet and helps to hold the upper layer together. Unless
you're a mechanical genius, you really can't understand this without a picture, or,
better yet, the real thing.
        In his definitive treatise, Notes on Rubik's `Magic Cube, David Singmaster,
professor of Mathematical Sciences and Computing at the Polytechnic of the South
Bank in London, defines the basic mechanical problem as that of figuring out how the
cube is constructed. I sometimes wonder whether Rubik's intended visualization task
for his students was to solve the unscrambling problem (Singmaster calls it the basic
mathematical problem) or to solve the mechanical problem. I suspect the latter is the
harder of the two. I myself must have put in more than 50 hours of work, distributed
over several months, before I solved the unscrambling problem, and I never did solve
the mechanical problem until I saw the cube disassembled. Singmaster informally
estimates that people who eventually solve the unscrambling problem (without hints)
take, on the average, two weeks of concentrated effort. Of course, it is hard for
anyone who has done it to say exactly how long it took (how can you tell play from
work?), but it's safe to say that if you are destined to solve the unscrambling problem
at all, it will take you somewhere between five hours and a year. I trust this is
reassuring.
        An important fact that many people fail to appreciate at first is that to restore a
scrambled cube even once to the START position (the state of Perfect Enlightenment
and Grace, where each face is a solid color) is so hard that it is necessary to find a
general algorithm for doing it from any scrambled state. No one can restore a messed-
up Magic Cube to its pristine state by




Magic Cubology                                                                         303

FIGURE 14-2. In (a) the three types of cubie are identified: face centers (F), corners
(C), and edges (E). In (b) the mechanism is revealed You can see the six pronged
infernal spindle with all six face-center cubies attached to it, and one detached edge
cubie,- and one detached comer cubie. Notice that no cubie,- is a complete cube. In
fact, the face centers are just façades! In (c), the gradual dismantling and rebuilding
of a Cube are shown. Warning: If you follow this procedure, you are advised to
rebuild your Cube in its pristine state,• otherwise, you will probably wind up with
your Cube in an orbit from which START is inaccessible




Magic Cubology                                                                     304

 mere trial and' error'. Anyone who gets back to START has built up a small science.
        A word of warning: Proposed solutions to the mechanical problem are often
lacking in clarity, having either too much or too little detail.- It is certainly a challenge
to come up with a mechanism that has the multifaceted twistability of the Magic
Cube, but it is perhaps no less of a challenge to describe the mechanism in language
and diagrams that other people can readily comprehend. By the same token, to convey
algorithms that restore the cube to START calls for a good, clear notation. Singmaster
himself has an excellent notation that is now considered standard; I will present it
below. A second word of warning: I am not a "cubemeister" (one who has contributed
to the annals of the profound science of Cubology); I am a mere cubist, an amateur
dazzled by the Cube and by the virtuosos who have mastered it. Therefore I am not a
suitable recipient of novel solutions to the mechanical problem or to the unscrambling
problem. I recommend to readers who believe that they have some novel insight to
communicate it to Singmaster, who runs what amounts to the World Center for
Cubology. lis address is: Department of Mathematical Sciences and Computing,
Polytechnic of the South Bank, London SEI OAA, England.

                                         *   *   *

        By now, I would hope that your appetite has been whetted to the point where
immediate possession of a Magic Cube is an urgent priority. Fortunately, this can be
arranged quite easily. Most any toy store now carries them under such names as
"Rubik's Cube", "Wonderful Puzzler", and miscellaneous others. The price ranges
from a couple of dollars for a cheap model to roughly $15 for a very solid and high-
quality cube. It is likely that many people will-buy cubes, little suspecting the
profound difficulty of the "basic mathematical problem". They will innocently turn
four or five faces, and suddenly find themselves hopelessly lost. Then, perhaps
frantically, they will begin turning face after face one way and then another, as it
dawns on them that they have irretrievably lost something precious. When this first
happened to me, it reminded me of how I felt as a small boy, when I accidentally let
go of a toy balloon and helplessly watched it drift irretrievably into the sky.
        It is a fact that the cube can be randomized with just a few turns. Let that be a
warning to the beginner. Many beginners try to clave their way back to START by
first getting a single face dope. Then, a bit stymied, they leave their partially solved
cube lying around where a friend may spot it. The well-known "Don't touch it!"
syndrome sets in when the friend innocently picks it up and says, "What's this?" The
would-be solver, terrified that all their hard-won progress will be destroyed, shrieks,
"Don't touch it!" Ironically, victory can come only through a more flexible attitude
allowing precisely that destruction.




Magic Cubology                                                                          305

        For the beginner, there is awesome sense of 'irreversibility about destroying
START, a fear of tumbling off the edge of a precipice. When my own first cube (1
now have dozens) was first messed up (by a guest), I felt both relieved (because it was
inevitable) and sad (because I feared START was gone forever). The physicist in me
was reminded of entropy. Once START had become irretrievable, each new twist of
ore face or another seemed irrelevant. To my naïve eye there, was no distinguishing
one messed-up state from another, just\_ as to the naïve eye there is no distinguishing
one plate of spaghetti from another, one pile of fall leaves from another, and so on.
The details meant nothing to me, so they didn't register. As I performed my "random
walk", the vastness of the space of possible shufllings of the little cubies became
vivid.
        As with a deck of cards, one can calculate the exact number of possible
rearrangements of the cube. An initial estimate would run this way. The first
observation-a rather elementary one-is that on the rotation of any face, each comer
goes to another comer, each edge to another edge and the center of the face stays put
(except for its invisible rotation). Therefore corners mix only with their own kind, and
the same goes for edges. There are eight comer cubies and eight comer cubicles (the
spatial niches, regardless of their content). Cubies and cubicles are to the cube as
children and chairs are to the game of musical chairs. Each comer cubie can be
maneuvered into any of the eight comer cubicles. This means that we have eight
possible fillers for cubicle No. 1, seven for cubicle No. 2, six for cubicle No. 3, and so
on. Therefore the corners can be placed in their cubicles in 8 X 7 X 6 X 5 X 4 X 3 X 2
X 1 (= 81) different ways. But each corner can be in any one of three orientations.
Thus one would expect a further factor of 38 from the eight corners. One would
expect the same for the twelve edge cubies: twelve objects can be permuted among
themselves in 12! different ways, and then, since each of them has-two possible
orientations, that gives another factor of 212. The center cubies never leave their
START positions (unless the cube is rotated as a whole) and have no visibly distinct
orientations, so they do not contribute. If we multiply the numbers out, we get
519,024,039,293,878,272,000 possible positions-about 5.2 X 1020.
        But there is an assumption here: that any cubie can be gotten into any cubicle
in any orientation, regardless of the other cubies'' positions and orientations. As we
will see, this is not quite the case. It turns out that there is a mild constraint on the
orientation of the comer cubies: any seven can be oriented arbitrarily, but the last one
is then forced, thus removing-one factor of three. Similarly, there is a mild constraint
on edge cubies: of the twelve, any eleven can be oriented arbitrarily, but the last one is
then determined, so that another factor of two is removed. There is one final
constraint on the permutations of cubies (disregarding their orientations) that says you
can place all but two of them wherever you want, but the last two are forced. This
removes a final factor of two, reducing the estimate above by a total factor of 3 X 2 X
2 =12, bringing the possibilities down to




Magic Cubology                                                                        306

a mere 48,252,003,271;489,856,000---aboul-4 3 X 1019 Still, `it must be said , this
does slightly exceed the assertion on Ideal's label: "Over three billion combinations".
        Another way of thinking about this factor of twelve is that if you begin at
START, you are limited to a twelfth of the "obvious" states, but if you disassemble
your cube and reassemble it with a single comer cubie twisted by 120 degrees, you are
now in a formerly inaccessible state, from which a whole family of
43,252,003,274,489,856,000 new states is accessible. There are twelve such
nonoverlapping families of states of the cube, usually called orbits by group theorists.

                                       *   *   *

         Speaking of impossible twists, I would like to mention a lovely discovery in
Cubology that is parallel to ideas in particle physics. It was pointed out by
mathematician Solomon W. Golomb. The discovery states: It is impossible to find a
sequence of moves that leaves just one comer cubie twisted a third of a full turn and
everything else the same. Now, recalling the famous hypothetical fundamental
particle with a charge of + 1/3 and its antiparticle with a charge of -1/3, Golomb calls
a clockwise one-third twist a quark and a counterclockwise- one-third twist an
antiquark. Like their cubical namesakes, quark particles have proved to be
tantalizingly elusive, and particle physicists generally believe now in quark
confinement: the notion that it is impossible to have an isolated free quark (or
antiquark). This correspondence between cubical quarks and particle quarks is a
lovely one.
         Actually, the connection runs even deeper. Although quark particles cannot
exist free, they can exist bound together in groups: a quark-antiquark pair is a meson
(Figure 14-9e), and a quark trio with integral charge is a baryon. (An example is the
proton-qqq-with a charge of + 1.) Now in the Magic Cube, amazingly enough, it is
possible to give any two comer cubies one-third twists, provided they are in opposite
directions (one clockwise, the other counterclockwise). It is also possible to give any
three comer cubies one-third twists, provided they are all in the same -direction. Thus
Golomb calls a- state with two oppositely twisted corners a "meson", and one with
three corners twisted in the same direction a "baryon". In the particle world, only
quark combinations with an integral amount of charge can exist. In the cubical world,
only quark combinations with a integral amount of twist are allowed. This is just
another way of saying that the orientation of the eighth comer cubie is always forced
by the first seven. In the cubical world, the underlying reason for "quark confinement"
lies in the group theory. There may be a closely related group=theoretical explanation
for the confinement of quark particles. That remains to be seen, but in any event, the
parallel is provocative and pleasing.

                                       *   *   *




Magic Cubology                                                                      307

         If we have a "pristine cube" (one=in START); what kind of move sequence
will create a meson or a baryon? Here we have an example of the most powerful idea
in Cubology: the idea of "canned" move sequences that accomplish some specific
reordering of a few cubies, leaving everything else untouched ("invariant", as group
theorists say). There are many different terms for such canned move sequences. I have
heard them called operators, transforms, words, tools, processes, maneuvers, routines,
subroutines, and macros, the first three being group-theoretical terms and the last
three being borrowed from computer science. Each term has its own flavor, and I find
that I use them all at various times.
In order to talk about processes, we need precision, and that means a good technical
notation. I will therefore present Singmaster's notation now. First we need a way of
referring to any particular face of the cube. One possibility is to use the names of
colors as the names of the faces, even after the cubies have become mixed up. Now it
might seem that calling a face "white" would be meaningless if white is scattered all
over the place. But remember that the white center cubie never moves with respect to
the five other center cubies, and thus defines the "home face" for white. So why not
use color names for faces? Well, one problem is that different cubes come with their
colors arranged differently. Even two cubes from one manufacturer may have
different START positions.. A more general convention is to refer to faces simply as
left and right, front and back, and top and bottom. Unfortunately, the initials of "back"
and "bottom" conflict. Singmaster resolves the conflict by replacing "top" and
"bottom" by up and down. Now we have names for the six faces: L, R, F, B, U, D.
Any particular cubie can be designated by lowercase italic letters naming the faces it
belongs to. Thus ur (or ru) stands for the edge cubie on the right side of the top layer,
and urf for the corner cubie in front of it (see Figure 14-3a).
The most natural move for a right-handed cubist seems to be to grasp the right face
with the thumb pointing up along the front face and to move the thumb forward. Seen
from the right side, this maneuver causes a clockwise quarter-twist of the R face. This
move will be designated R (see Figure 14-3b)
.
The mirror-image move, where the left hand turns the L side counterclockwise (as
seen from the left), is L-t, or, for short, L'. A clockwise twist of the L side is called,
naturally, L. A 90-degree clockwise turn of any face (from the point of view of an
observer looking at the center of that face) is named by the letter for that face, and its
inverse-the counterclockwise quarter turn --has a prime mark following the face's
initial. Quarter-turns will henceforth be called q-turns.
With this nomenclature, we can now write down any move sequence, -no matter how
complex. A trivial example is four successive R's, which we write as R4. In the
language of group theory, this is the identity operation: it has zero effect. An equation
expressing this fact is R4 = I. Here, I stands for the "action" of doing nothing at all.
Suppose we twist two different faces-say R first, then U. We will




Magic Cubology                                                                        308

FIGURE 14-3. Labeling of rubies and moves. The speckled cubie in (a) is the urf
cubie (alias rfu and fur), and the black one is the ur (or ru) cubie. The quarter-turn or
q-turn shown in (b) is called R.




Magic Cubology                                                                       309

transcribe that as RU—nor as UR Note; in fact, that RU and UR are quite different in
their effects. To check this out, first perform RU on a pristine cube, observe its
effects, then undo it, try UR, and see how its effects differ: The inverse of RU is, quite
obviously, U'R', not R'U'. (Incidentally, this strategy of experimenting with move
sequences on a pristine cube is most helpful. Very early I found it useful to buy a
second cube so that I could work on solving one while experimenting with the other,
never letting the second one get far away from START)

                                        *   *   *

         What is the effect of a particular "word"? That is to say, which cubies move
where? To answer this question, we need a notation for the motions of individual
cubies. The effect of R on edges is to carry the ur cubie around to the back face to
occupy the br cubicle. At the same time, the br cubie swings around underneath,
landing in the dr position, the dr cubie moves up like a car on a Ferris, wheel to fill
theft cubicle, and theft cubie comes to the top at ur. (See Figure 14-4a.) This is called
a 4-cycle, and we'll write it in a more compact way: (urbrdrfr). Of course, it does not
matter where we start writing; we could equally well write (brdrfr,ur).
         On the other hand, the order of the letters in cubie names does matter. We can
reverse all of them or none of them, but not just some of them. If you think of the
letters as designating facelets, this will become clear. For example, if we wrote
(urrb,dr,rf), it would represent a 4-cycle involving the same four cubicles as above,
"but one in which each cubie flipped before moving from one cubicle to the next. Of
course, such a cycle cannot be accomplished by a single q-turn, but it may be the
result of a sequence of q-turns of different faces (an operator). Or consider the
following 8-cycle, shown in Figure 14-4c: (uruf ul,ub,rufu,lu,bu). This has length
eight, but involves only four cubicles. Each cubie, after making a full swing around
the top face, comes back flipped (see Figure 14-4b). After two full swings, it is back
as it started. Each facelet has made a "Mobius trip". We can designate this "flipped 4-
cycle" as '(uruf ul\_ub)+, where the plus sign designates the flipping. The designation
(ru,fu,lu,bu)+ and numerous others would do as well. Thus the cycle notation tells you
not only where a cubie moves but also its orientation with respect to the other cubies
in its cycle.
         To complete our description of the effect of R, we must transcribe the 4-cycle
of the corners. As with edges, we have the freedom to start at any corner we want, and
once again we must be careful to keep track of the facelets so that we get the
orientations right. Still, R has a rather-trivial effect on corners: (urf bru,drb,frd),
which could also be written (rub,rbd,rdf rfu), and many other ways. Summing up, we
can write R = (urbrdr r) (urf,bru,drb,frd). This says that R consists of two disjoint 4-
cycles. (If we wanted to, we could throw in a term standing for the 90-degree rotation
of the R face's center, but since such rotation is invisible, we needn't do so.)




Magic Cubology                                                                        310

FIGURE 14-4. The simple 4-cycle (ur,br,dr,rr), shown in (a), is what happens to edge
cubies during the q.turn R. In (b), a trickier 4-cycle (ur,rb,dr,rf), involving the same
four cubies, is shown; here, each cubie flips before entering the next cubicle. This
cycle can be produced only through a sequence of q-turns. In (c), the 8-cycle (fi-
,ur,br,dr,rf ru,rb,rd) is shown, which can also be thought of as a 'flipped 4-cycle'-
namely, (fr,ur,br,dr)+. In (d), the 7-cycle (ur,br,dr, fr,uf,ul,ub) is shown snaking its
way around the Cube, representing the. effect on edges of the simple operator RU.

        What about transcribing a move sequence such as RU? Well, take a pristine
cube and perform RU. Then start with some arbitrary cubie that has moved and
describe its trajectory. For example, ur has moved to br. Therefore br has been
displaced. Where has it gone? Find the new location of that cubie (it is dr) and
continue chasing cubies 'round and 'round the cube until you find the one that moved
intoo the original position of ur. You will find the following 7-cycle: (ur,brdr fr,uf
ul,ub) (see Figure 14-4d).
        What about corners? Well, suppose we trace the cubie that originated in




Magic Cubology                                                                      311

where did RU carry it? The answer is: Nowhere- it took a round trip but got twisted
along the way. It changed into rfu. We can designate this clockwise twist-this "twisted
unicycle", this quark-as (urf)+. This is shorthand for the following 3-cycle: (urf,rju
fur). You can even see this as cycling the three letters u, r, and f inside the cubie's
name. If the cycle had been an antiquark, we would have written (urf) -, and the
letters would cycle the other way.
        What about the other seven corners? Two of them-dbl and d4(-stay put, . and
the other five almost form a 5-cycle: (ubrbdrdfr,luf,bul). It is unfortunate that the
cycle does not quite close, because bul, although it gets carried into the original ubr
cubicle, does so in a twisted manner. It gets carried to rub, which is a
counterclockwise twist away from ubr. This means we are dealing with a 15-cycle.
But it is so close to the 5-cycle above that we'll just tack on a minus sign to represent
the counterclockwise twist. Our twisted 5-cycle is then (ubr,bdr,dfrluf bul)\_, and the
entire effect of RU, expressed in cycle notation, is (ur,brdrfr,uf,ul,ub) (urf)+
(ubrbdrdfr,luf,bul)\_.
Now that we have RU in cycle notation, we can perform rotations mentally, by sheer
calculation. For instance, what would be the effect of (RU)s"? Edge cubie ur would be
carried five steps forward along its cycle, which would bring it to ul. (This can also be
seen as moving two steps backward.) Then ul would go to fr and so on. The 7-cycle is
replaced by a new 7-cycle: (ur,ul fr,brub,uf dr). Let us now look at the twisted 5-
cycle. Corner cubie ubr would be carried five steps forward along its cycle, which
brings it back to itself negatively twisted-namely, rub. Similarly, all the corner cubies
in the 5-cycle would return to their starting points, but negatively twisted; thus, on
being raised to the fifth power, a negatively twisted 5-cycle becomes five antiquarks.
But if that is so, . how is the requirement for integral twist satisfied? Don't we have
one quark-(urf)+and five antiquarks, and doesn't that add up to four antiquarks, with a
total twist of -1 3? Well, I have slipped something by you here. Can you spot it?
o gain facility with the cycle notation, you might try to find the cycle representation of
various powers of RU and UR and their inverses.

                                        *   *   *

Any sequence of moves can be represented in terms of disjoint cycles of various
lengths (cycles with no common elements). If you are willing to let cycles share
members, however, any cycle can be further broken up into -cycles (called
transpositions, or sometimes swaps). For instance, consider three animals: an
Alligator, a Bobcat, and a Camel. They initially occupy three ecological niches: A, B,
and C (see Figure 14-5). The effect of the




Magic Cubology                                                                        312

FIGURE 14=5. A zoological 3-cycle involving three objects: a, b, and c (an alligator,
a bobcat, and a camel). Initially, each is in its usual ecological niche: a in A, b in B,
and c in C. But then, after a permutation, c is in A, a is in B, and b is in C. This 3-
cycle can bethought of as the result of two successive swaps.




Magic Cubology                                                                       313

3-cycle, (A;B,C) is to put them' in the order Camel, Alligator, Bobcat. 'The same
effect can be achieved, however, by first performing the swap (A,B) (what was in A
goes to B and vice versa) and then performing (A, C). Of course, this can also be
achieved by the two successive swaps (A,C)(B,C)or, for that matter, by (B,C)(A,B).
On the other hand, no sequence of three swaps will achieve the same effect as
(A,B,C). Try it yourself and see. (Note ,that a niche is like a cubicle and an animal is
like a cubie.)
        An elementary theorem of zoop theory (a field we won't go into here) states
that no matter how a given permutation of animals among niches is reduced to a
product of successive swaps (which can always be done), the parity of the number of
such swaps is invariant; that is to say, a permutation cannot be expressed as an even
number of swaps one time and an odd number another time. Moreover, the parity of
any permutation is the sum of the parities of any permutations into which it broken up
(using the rules for addition of even and odd numbers: odd plus even is odd, and so
forth.)
        Now, this theorem has repercussions for the Magic Cube. In particular, you
can see that any q- turn consists of two disjoint 4-cycles (one on edges and one on
corners). What is the parity of a 4-cycle? It is odd, as you can work out for yourself.
Thus, after' one q- turn, both the edges and the corners 'have been permuted oddly;
after two q- turns, evenly; after three q-turns, oddly; and so forth. The edges and
corners stay in phase, in the sense that the parities of their permutations are identical.
Now clearly, the null permutation is even (it effects zero swaps). So if we have a null
permutation on corners, the permutation on edges must also be even. Conversely, a
null permutation on edges implies an even permutation on corners. Imagine a state
identical to START except for two interchanged edges (that is, one swap).. Such a
state would be even in corners but odd in edges, hence impossible. The best we could
do would be to have two pairs of interchanged edges. The same argument holds for
corners. In 'short, we have proven that single swaps are impossible; swaps must
always come in pairs. (This is the origin of one of those factors of two in the earlier
calculation of the number of reachable states of the cube.) There are processes for
exchanging two pairs of edges, two pairs of corners, and even for exchanging one pair
of edges along with one pair of corners. (This last process necessarily involves an odd
umber of q- turns.)
        To round out the subject of constraints, let us ponder the origin of the
constraints on corner-twisting and edge-flipping. Here is a clever explanation
provided by John Conway, Elwyn Berlekamp, and Richard Guy, elaborating an idea
due to Anne. Scott. The basic concept is that we want to show that the number of
flipped cubies is always even, and that the twist is always integral. But in order to
determine what is flipped and what is twisted, we need a frame of reference. To
supply it, we will define two notions: the chief facelet of a cubicle and the chief color
of a cubie. (Remember that a cubicle is a niche and a cubie is a solid object.) The
chief facelet of a cubicle will be the one on the up or down surface of the cube, if that
cubicle




Magic Cubology                                                                        314

FIGURE 14-6. Diagrams to aid in the proof that flippancy is even and twist is
integral.
        In (a), the Cube is in START. The chief facelets of cubicles are shown by
crosses and the chief colors of cubies by circles. (Note: The concept of "chiefness"
does not apply to face-center cubicles or rubies.) Think of the crosses as floating in
space and the circles as being attached to the Cube,. so that when turns are made, the
crosses stay where they are but the circles move. The bottom face looks identical to
the top, the left face identical to the right, and the back face identical to the front.
        In (b), the results of the q-turn F are shown. The two empty circles indicate
that the two rubies they are attached to have lost their "sanity" : For them to regain
their sanity, one cubie would have to be twisted one-third clockwise while the other
cubie was twisted one-third counterclockwise, thus canceling each other's
contribution to the total twist of the Cube. Similar remarks apply to the invisible left-
hand face.
        In (c), the results of the q-turn R (as applied to START) are shown. Empty
circles again come in pairs. The top and bottom corner rubies on the front face (each
with an empty circle) have canceling twists, as in (b). The top and bottom edge rubies
on the tight face have canceling flippancies, and the one seemingly unmatched empty
circle (on the edge cubic on the front face) is paired with an empty circle on the
invisible back face.

is one; otherwise it will be the one on the left or right wall (see Figure 14-6). There
are nine chief facelets on U, nine on D, and four on the equator. (We - can ignore the
centers, because they never can be flipped or twisted.) The chief color of a cubic is
defined as the color that should be on the cubie's chief facelet when the cubie "comes
home" to its proper cubicle in the START position.
        Now the argument goes this way. Suppose the cube is scrambled. Any, cubie
that has its chief color in the chief facelet of its current cubicle will be called sane;
otherwise it will be called flipped (this applies to edge rubies) or twisted (this to
corner cubies). Obviously, there are two ways a cubie can, twisted: clockwise (+ 1/3
twist) and counterclockwise (-1/3 twist). The flippancy of a cube state will be defined
as the number of flipped edge cubies in it, and the twist as the sum of the twists of the
eight corner rubies We shall say that the flippancy and twist of START are both zero,
by convention.
        Next consider the twelve possible q- turns out of which everything else is
compounded. Performing U or D (or their inverses) preserves both the


Magic Cubology                                                                       315

flippancy and the twist, since nothing leaves or enters the up or down face.
'Performing F or B (or their inverses) leaves the total twist constant, by changing the
twist of four corners at once: two by + 1/3 and two by -1/S'. It also leaves the
flippancy alone (see Figure 14-6b). Performing L or R will likewise leave the total
twist constant (four corner twists again cancel in pairs) and will change the flippancy
by 4, since always four cubies will .change in flippancy (see Figure 14-6c). The
conclusion is what I stated above without proof: the eight corner cubies are always
oriented to make the total twist a whole number, and the twelve edge cubies must
always be oriented to make the total flippancy even.

                                       *   *   *

       After this discussion of constraints, you should be convinced that no matter
how you twist and turn your Magic Cube, you cannot reach more than a twelfth of the
conceivable "universe", beginning at START It is another matter, though, to show
that every state within that one-twelfth universe is accessible from START (or what
amounts to the same thing, only backward: that START is accessible from every state
in the one-twelfth universe). For this, we need to show how to achieve all even
permutations of cubies, and how to achieve all orientations that do not violate the two
constraints described above. What it comes down to is that we have to show there are
operators that will perform seven classes of operations:

       (1) an arbitrary double edge-pair swap,
       (2) an arbitrary double corner-pair swap,
       (3) an arbitrary two-edge flip,
       (4) an arbitrary meson,
       (5) an arbitrary 3-cycle of edges,
       (6) an arbitrary 3-cycle of corners, and
       (7) an arbitrary baryon.

        Of course, each of these operators should work -without causing side effects
on any other parts of the cube. With these powerful tools in our kit, we would be able
to cover the one-twelfth universe without any trouble. In the case of the overlapping
swaps of animals, you saw how a 3-cycle is really two overlapping 2-cycles. This
implies that classes 5 and 6 can be made out 14 the first four classes. Similarly, a
baryon can be made from two overlapping mesons. -So all we really need is the first
four classes.
        To show that all the operators belonging to these four classes are available,
we'll use another of the most crucial and lovely ideas of Cubology: that of conjugate
elements. It turns out that all we need is one example in each class; given one
example, we can construct all the other operators of its class from it. How does this
work? The idea is very simple.
        Suppose we had found one operator in class 1 that swapped, say, of with




Magic Cubology                                                                     316

FIGURE 14-7. How to use conjugate moves to turn an unsolved problem into a solved
one. The unsolved problem is to effect the double swap shown by the white arrows.
The solved problem is the double swap shown by the black arrows (on top). As long
as we can maneuver the black cubies into the white cubicles, we are home free. This
principle has nothing to do with the specific cubicles involved in the known and
unknown operators, but simply with the idea that sometimes you can translate an
unsolved situation into a solvable situation, use a known operator to handle that
situation, then "back-translate " to regain the original situation,, but with the tricky
part now solved. This is the principle of conjugates.

ub, and ul with ur, leaving the rest of the cube undisturbed. Let us call this operator H.
Now suppose we wanted to swap two totally different pairs of edge cubies, say ft with
fd, and rb with rd (see Figure 14-7). We can daydream: "If only those cubies were in
the four `magical swapping spots' on the top surface..." Well, why not just put them
up there? It would be fairly simple to get four cubies into four specific cubicles. The
obvious objection is: "Yes, but that would have an awful side effect-it would totally
mess up the rest of the cube." But there is a clever retort. Let the destructive maneuver
that gets those four cubies into the magical swapping spots be called A. Suppose we
were smart enough to transcribe the move sequence of A. Then right after performing
A; we perform our double swap H. Now comes the clever part. Reading our transcript
in reverse order and inverting each q-turn, we perform the exact inverse of A. This
will not only un-maneuver the four cubies back into their old cubicles, but will also
undo the side effects A created in the rest of the cube. Does that restore the cube
intact? Not quite. Remarkably, since we sandwiched H between A and A', the four
edge cubies go home permuted-that is, each one winds up in the home of its swapping
partner! Other than that, the cube is restored,' and so we have accomplished precisely
the double swap we set out to accomplish.
        When you think this through, you see that it is flawless in conception. The
inverse maneuver, A', does not "know" we have exchanged two pairs of edges. As far
as it is concerned, it is merely putting -everything back where it was before A was
executed. Hence we have "snuck" our swaps in under A"s nose, which is to say we
have "fooled the cube". Symbolically, we have




Magic Cubology                                                                        317

carried out the sequence of moves AHA', which is called a conjugate of H.
        It is this kind of marvelously concrete illustration of an abstract notion of
group theory that makes the Magic Cube one of the most amazing things ever
invented for teaching mathematical ideas. Normally, the examples of conjugate
elements given in group-theory courses are either too trivial or too abstract to be
enlightening or exciting. The Magic Cube, though, provides a vivid illustration of
conjugate elements and of many other important concepts of group theory.

                                      *   *   *

         Suppose you wanted to get a quark-antiquark pair on opposite corners, but
knew how to do so only on adjacent corners. How could you do it? Here is a hint:
There are two nice solutions, but the shorter and prettier one involves using a
conjugate. Incidentally, any maneuver that creates a quark on one corner (with other
side effects, of course) might be called a quarks crew.
         What we have shown for edges goes also for corners: the ability to swap two
specific corners enables you to swap any two corners. Conjugation allows you to
build up an entire class of operators from any single member of that class. Of course,
the question still remains: How do you find some sample operator in each of the four
classes? For example, how do you find an operator that creates a meson on two
adjacent corners (a combination of a quarkscrew and an antiquarkscrew)? How do
you find an operator that exchanges two edge pairs both of which are on the top
surface? I won't give the answer here, but will follow Singmaster, who points the way
by suggesting quasi-systematic exploration of some small "subuniverses within the
totality of all cube states-that is, he suggests you look at subgroups. This means
restricting your set of moves deliberately to some special types of move. Here are a
few examples of interesting subgroups created by various kinds of restriction:

          1. The Slice Group. In this subgroup, every turn of one face must be
  accompanied by the parallel move on the opposing face. Thus RR must be
  accompanied by L', U by D', and F by B'. The name comes from the fact that
  any such double move is equivalent to rotating one of the three central slices of
  the cube. Singmaster abbreviates the slice move RL' by R5, R'L by Rs, and so
  forth, Under this restriction, faces cannot get arbitrarily scrambled. Each face
  will have a pattern in which all four corners share one color (Figure 14-8). A
  special- case is the pattern called Dots, in which each face is all one color
  except for its center (see Figure 14-9a). Can you figure out how to achieve Dots
  from START? How many different ways are there of arranging the -dots? How
  does the Dots pattern resemble a meson? (You will find answers to all these
  questions, along with much else, in Singmaster's book.)




Magic Cubology                                                                    318

       FIGURE 14-8. The type of pattern that the Slice Group creates on all faces.

           2. The Slice-Squared Group. Here we restrict the Slice Group further,
   allowing only squares of slice moves, such as R52 (which is the same as R2L)
   or F52 (which is the same as F2B2).
           3. The Antislice Group. Here, instead of always rotating opposing faces
   in parallel, we always rotate them in antiparallel, so that R is accompanied by L,
   F by B, and U by D. An antislice move has a subscript a, as in R., which equals
   RL. (Of course, the Antislice-Squared Group is no different from the Slice-
   Squared Group.)
           4. The Two-Faces Group. Allow yourself to rotate only two adjacent
   faces, say F and R. It turns out to be a pretty substantial challenge to figure out
   an algorithm for undoing an arbitrary scrambling of two faces, staying within
   the Two-Faces Group. Most cube experts will instead resort to the "elephant
   gun" of twisting all six faces to get out of a mere two-face scramble. Shame on
   them!
           5. The Three-Faces Groups. The reason this category is pluralized is that
   there are nonequivalent choices of threesomes of faces. For example, you can
   form a kind of "bridge", as with faces L, U, and R, or you can form a "corner",
   as with faces F, U, and R.
           6. The Four-Faces and Five-Faces Groups. Again, there are various
   non-equivalent choices of faces for four faces. The Five-Faces Group is, as it
   turns out, actually the full group of the cube. In other words, you can make an
   operator equivalent to R out of L, U, D, F, and B.
           7. -The Two-Squares Group. As in the Two-Faces Group, you may
   rotate only two faces, using only 180-degree turns at that. This is a very simple
   subgroup.

        If you limit your attention to just the Two-Faces and Two-Squares groups, you
will be able to find processes that achieve double swaps-some of edges, others of
corners. It is a remarkable fact that these processes alone, together with the notion of
conjugation, will allow us-in a theoretical sense-to solve the entire unscrambling
puzzle.
        Why don't we also need a meson maker and, a double edge-flipper? Well,


Magic Cubology                                                                       319

consider how we might make a double edge flipper from the two classes of tools one
may assume will be found-that is, double edge-swappers and double corner-swappers.
In order to flip two edges without creating any side effects, we'll perform two
successive double edge-pair swaps, and both times they will involve the same pairs!
For example, we might swap of with ub, and df with db, and then reswap them. This
seems to be an absolute "nothing process", but that need not be the -case. After all,
just as before, we can sandwich the second swap between a process X and its inverse
X', where we carefully choose the process X so as to... (Oh, darn it all, I totally lost
my train of thought there. I'm sure you can finish it up, though. I do remember that it
wasn't too tricky, and that I thought the idea was rather elegant. I'm sure you will too.)
        The same kind of thinking will show how you can build up a meson maker out
of mere corner-swapping processes and conjugation. Given mesons, you can build up
baryons. And with mesons and baryons, double edge-flippers, double edge-pair
swappers, and double corner-pair swappers, you have a full kit of tools with which to
restore any scrambled cube to START, as long as it belongs to the same orbit as
START. What I have given is, needless to say, a highly theoretical existence proof,
and any practical set of routines would be organized quite differently. The type of
solution I have described has the advantage of being compact in description, but it is
enormously inefficient.-In practice, a cube solver must develop a fairly large and
versatile set of routines that are short, easy to memorize, and highly redundant. There
is an advantage to being able to carry out transformations in a variety of ways: you
can choose whichever tool seems best adapted to the situation at hand, instead of, for
instance, using some theoretically developed tool that takes several hundred q-turns to
make a baryon.

                                           * * *
         The typical cube solver evolves a set of transforms partly by intuition, partly
by luck, sometimes with the aid of diagrams, and occasionally with abstract principles
of group theory. One principle nearly everyone formulates quite early is that of
"getting things out of the way". This is once again the idea of conjugates, only in a
simpler guise. The typical patter that goes along with it is something like this (I have
included sound effects of -a sort): "Let's see, I'll swing this out of the way [flip, flip]
so that I can move that [flap, flap], and now I can swing this back again [unflip,
unflip]. There -now I've got that where I wanted it to be." You can hear the conjugate
structure inside the patter ("flip, flap, unflip").
         The only problem with being conscious of why it all works as you carry it out
is that it may be too taxing. My impression is that most cubemeisters do not think in
much detail about how their tools are achieving their goals, at least not while they are
in the midst of restoring some scrambled cube. Rather, expert cube solvers are like
piano virtuosos who have memorized




Magic Cubology                                                                         320

difficult pieces. As Dan' Weise, an -MIT cubemeister, said to me, e forgotten how to
solve the cube, but luckily, my fingers remember."
        The average operator seems to be about ten to twenty q- turns long. You don't
ever want to get lost in mid-operator, because if you do, you will have a totally'
scrambled cube on your hands, even if you were carrying out your final transform. As
cubemeister Bernie Greenberg said to me once, "If I were solving a cube and
somebody yelled `Fire!', I would finish my transform before clearing out."
        My own style is probably overly blind. Not only do I not think about why my
operators work as I am carrying them out; I have to admit that with some of them, I
don't even have the foggiest idea why they work at all! I found these "magic
operators" through a long and arduous trial-and-error procedure. I used some heuristic
notions, such as: "Explore various powers of simple sequences", "Use conjugates a
lot", and so on. One thing I hardly used at all-alas, poor Rubik-was three-dimensional
visualization. However, I do know one Stanford cubemeister, Jim McDonald, who
can give the reason for every last q-turn he makes. His operators don't seem magical
to him because he can see what they are doing at every moment along the way. In
fact, he does not have them memorized as I do mine; he seems to reconstruct them as
he unscrambles cubes, relying on his "cube sense". He is like an expert musician who
can improvise where a novice must memorize. For interested readers, the central idea
of Jim's method is first to solve the top layer except for one corner, and then to utilize
the vertical "chimney" underneath that free corner as you might use a neighbor's
driveway to turn your car around in. The other two layers are cleaned up by shunting
cubies in and out of the "chimney/driveway".

                                        *   *   *

        Perhaps not coincidentally, the abstract approach has been carried -to its
extreme by Singmaster's officemate, Morwen B. T.histlethwaite (I wonder what that
"B" stands for!). He currently holds the world record for the shortest unscrambling
algorithm. It requires at most 52 "turns". (A turn is defined as: either a q- turn or a
half-turn-that is, a 180-degree turn of one face.) Thistlethwaite has used ideas of
group theory to guide a computer search for special kinds of transforms. His
algorithm has the curious property of not giving any appearance of converging toward
the solved state at all-until the very last few turns.
        This must be contrasted with the - more conventional style. Most algorithms
begin by getting one layer-usually the top layer-entirely correct. (In saying "top layer"
rather than "top surface", I mean that the "fringe" has to be right, too: that is, the
cubies on top must be correct-as seen from the side as well as from above.) This
represents the first in a series of "plateau states". Although further progress requires
any plateau state's destruction, that state will later be restored, and each time this
happens,




Magic Cubology                                                                        321

more order will have been introduced. These are the successive plateau states.
        After getting the top layer, the solver typically works on corners on the bottom
layer, or perhaps on getting the horizontal equator slice all fixed up. Most algorithms
can, in fact, be broken up into five or six natural stages, corresponding to natural
classes of cubies that get returned to their home cubicles. My personal algorithm, for
instance, goes through the following five stages:

       (1) top edges,
       (2) top corners,
       (3) bottom corners,
       (4) equator edges, and
       (5) bottom edges.

In the first two of my stages, placement and orientation are achieved simultaneously.
Each of the last three stages breaks up into substages: a placement phase and then an
orientation phase. Naturally, the operators of any stage must respect all the
accomplishments of preceding stages. This means that they may damage the order
built up as long as they then repair it. They are welcome, however, to indiscriminately
jumble up cubies scheduled to be dealt with in later stages. I find that other people's
algorithms are usually based on the same classes of cubies, but the order of the stages
can be completely different.
        Virtually all algorithms have the property that if you were to take a series of
snapshots of the cube at the plateau states, you would see whole groups, of cubies
falling into place in patterns. This is called "monotonicity at the operator level"-that
is, a steady, visible approach toward START, with no backtracking. Of course, you
would see something totally different if you took snapshots between plateau states-but
that is another matter.. There is no known algorithm that makes visible progress with
every turn!
        Very different in spirit is Thistlethwaite's algorithm. Instead of trying to put
particular classes of cubies into their cubicles, he makes a "descent through nested
subgroups". This means that, starting with total freedom of movement, he makes a
few moves, then clamps down on the types of move that will thenceforward be
allowed, makes a few more moves, clamps down a bit more, and so on, until the
constraints become so heavy that nothing can move any more. But just at this point,
the START position has been achieved! Each time, the clamping-down amounts to
forbidding q- turns on two opposite faces, allowing only half-turns in their stead from
then on. The first faces to be thus "clamped" are U and D, then come F and B, and
finally L and R. The strange thing about this approach is that you cannot see START
getting nearer, even if you take a series of snapshots at carefully chosen moments.
Just all of a sudden, there it is! It's as if you were climbing Everest and the peak were
shrouded in clouds until the last 100 met .when suddenly the clouds break and there
it- is!




Magic Cubology                                                                       322

        This Thistlethwaite' algorithm thuggests a thorny thought: Wouldn't it be nice
if there were an easy way to tell how far you are from START? We might call this a
"distance-from-START-ometer". .Such a device would obviously be quite useful. For
example, it is rather embarrassing to resort to the full power of a general
unscrambling algorithm to undo what some friend has done with four or five casual
twists. For that reason alone, it would be nice to be able to assess quickly if some state
is "really random" or is close to START. But what does "close" mean? Distances
between two states in this vast space can be measured in two fairly natural ways. You
can count either the number of q-turns or the number of turns needed to get from one
state to the other (where "turn", as above, means either a q- turn or a half-turn). But
how can one figure out how many turns are needed to get to START without doing an
exhaustive search? A reliable and at least fairly accurate estimate would be preferable,
one that could be carried out quickly during a cursory inspection of the cube state. A
nave suggestion is to count the number of cubies that are not in\_ their home cubicle.
This estimator, however, can be totally fooled by the Dots position, in which nearly
all cubies are on the "wrong" side (see Figure 14-9a). That position is only eight q-
turns away from START. Perhaps the flippancy and the number of quarks could also
be taken into account by a better estimator, but I don't know of any.
        There are sophisticated group-theoretical arguments suggesting that-the
farthest one can get from START is 22 or 23 turns. This is quite striking, considering
that most solvers' early algorithms take several hundred turns, and highly polished
algorithms take a number somewhere in the 80's or 90's. Indeed, many mere operators
take considerably more turns than Thistlethwaite's entire algorithm does. (My first
double edge-flipper, for instance, was nearly 60 turns long.)
        One result that can be demonstrated easily is that there exist states at least 17
turns away from START. The argument goes as follows. At the outset there are 18
possible turns we might make: L, L', L2, R, R', R2, and so on. After that, there are 15
reasonable turns to make. (One would not move the same face again.) The number of
distinct turn sequences of length 2 is therefore 18 X 15, or 270. Another turn will
contribute another factor of 15;: and so on. How long does it take before we have
reached the number of accessible states? It turns out that 17 is the smallest number of
turns that, will theoretically allow access to 4.3 X 1019 distinct states. Of course, not
every turn sequence of length 17 leads to a unique state, not by a long shot, and so we
haven't shown that 17 turns will reach every accessible state. We have simply shown
that at least 17 turns are needed if you want to reach every state from START. So,
conceivably, no two states are much more than 17 turns away from each other. But
which 17 turns? That is the question.
        So far, only God knows how to get from one state of the Magic Cube to
another in the minimum number of turns. "God's algorithm" is, by definition, the
speediest recipe for solving the Cube from any state. A burning question of Cubology
is: Is God's algorithm just a gigantic table without any




Magic Cubology                                                                        323

FIGURE 14-9. A number of special configurations deserving of names. In (a), the
pattern known as Dots. In (b), Pons Asinorum. In (c), the Christman Cross. In (d), the
Plummer Cross. In (e), a Meson (showing what appears to be an isolated quark, but it
is actually\_ balanced by an antiquark on, the'opposiL corner). In (f), a Giant Meson,
consisting of a "giant quark" and a "giant antiquark" on opposite corners.




Magic Cubology                                                                    324

pattern in it, or is there a significant mount of pattern to it, i6,ifiiit elegant and short
algorithm based on it could be mastered by a mere mortal? Notice that possession of a
distance-from-START-ometer would be tantamount to possession of God's algorithm.
Given any scrambled -state, you tentatively try out all eighteen possible twists and
then choose one that brings you closer to START (Why must there always be one?)
Make it, and then repeat the process. It's a little arduous, but it gets you to START
directly, obviating plateaus or other intuitive intermediary states. That's one reason for
doubting that any simple such meter exists.

                                         *   *   *

If God were to enter a cube-solving contest, It might encounter some rather stiff
competition from a few prodigious mortals, even if they do not know Its algorithm.
There is a young Englander from Nottingham named Nicholas Hammond who has got
his average solving time down to close to 30 seconds! Such a phenomenal
performance calls for several skills. The first is a deep understanding of the cube. The
second is an extremely polished set of operators. The third is to have the operators
down so cold that you could do them in your sleep. The fourth is sheer speed at
executing twisty hand motions. The fifth is having a well-oiled "racing cube": one that
turns at the merest twitch of a finger, eagerly anticipating every operator before it is
needed. In short, the racing cube is a cube that wants to win.
        I have not yet heard of people naming their racing cubes, although that is sure
to come. It would seem, though, that there is an correlation between having a colorful
name and being a contributor to Cubology. Apart from Singmaster and Thistlethwaite,
there is Dame Kathleen Ollerenshaw (late Lord Mayor of Manchester), who has
discovered many streamlined processes, has written an article on the Magic Cube, and
has the distinction of being the first to report an attack of Cubist's Thumb, a grave
form of the disease mentioned at the beginning of this column. Then there is Oliver
Pretzel, the discoverer of a delicious twisted 3-cycle and the creator of a lovely "pretty
pattern" called the "6-U" state, which can be reached from START by way of the long
word

       L'R2F'L'B' UBLFRU'RLR S F S U S R S .

        Pretty patterns are of interest to, many cube lovers, but I cannot do them
justice here. I can mention only a few of the best I know. A good warm-up. exercise is
to figure out how to make the state called Pons Asinorum ("Bridge of Asses"). It is
shown in Figure 14-9b. It has this name because, because as one MIT cubemeister
remarked to me, "If you can't hack this one, forget about cubing." Then there are two
kinds of cross, known to the MIT cube-hacking community as the Christman Cross
and the Plummer Cross (see parts (c) and (d) of Figure 14-9). The former involves
three pairs of colors




Magic Cubology                                                                         325

U-D, F-R, and' -L-B), while the latter involves two triples in the quark-antiquark
style. My favorite pretty pattern is the "Worm", whose "genotype", or turn sequence,
is:
        RUF2D'R S F S D'F'R'F2RU2FR2F'R'U'F'U2FR.

'Then there is the Snake, a similar sinuous pattern that winds around the cube:

       BR S D'R2DR',B'R2UB2U'DR2D.

If you cut off the Snake's tail (R2D') and instead stick on B2RaU2RaB2D', you will
create a curious bi-ringed pattern. All of these are from pretty-pattern-meister Richard
Walker. A beautiful pattern is the Giant Meson (Figure-14-9f), made from a giant
quark (a 2 X 2 X 2 corner subcube rotated 120 degrees) and a giant antiquark. If you
wish, you can top it off, using quarkscrews to twist a standard-size quark and
antiquark onto the corners of the giant quark and antiquark, like cherries on top of
sundaes. I'll let you figure out how to make this one.

                                       *   *   *

        I would like to leave you with a set of hints and some things to think about. A
difficult challenge, good for cubists at all levels of cubistry, is for someone to do a
handful of turns on a pristine cube, to return it to you in this mildly scrambled state,
and for you to try to get it back to START by finding the exact inverse word.
Cubemeisters will be able to invert a bigger handful of turns than novices. Kate Fried
reportedly can invert seven turns regularly, and once, after a full day of staring at the
cube, she undid ten. (I can undo about four.)
        My royal road to discovering an algorithm is based on two challenging
exercises involving corner cubies only. The preliminary exercise is as follows.
Maneuver the four corner cubies with white on them to the top face with- their white
facelets pointing upward. Do not worry about which cubie is in which cubicle.
Simultaneously do the same thing on the bottom face (of course with its color
pointing downward). The advanced exercise is to do the preceding one while in
addition making sure that all the corner cubies end up in their proper cubicles. This
amounts to solving the 2 X 2 X 2 Magic Cube puzzle, and it will take you a long way
toward mastery of the Magic Cube.
        To help you with your edge processes, here is a wonderful trick discovered by
David Seal, based on a type of operator called a monoflip. I'll give it to you as a
puzzle. How can you make a double edge-flipper out of a process that messes up the
lower two layers but leaves the top layer invariant, except for flipping a single edge
cubie? Hint: The answer involves the important group-theoretical idea of a
commutator-a word of the form




Magic Cubology                                                                       326

PQP'Q'. I will also leave it to you to find your own monoflip operator. After I found
out about it, I incorporated this trick into my method.
         . Here is a small riddle: Why do 5- and 7- cycles crop up so often in an object
whose symmetries all have to do with numbers such as 3, 4, 6, and 8? Where do cycle
lengths such as 5 and 7 come from? A somewhat related question is: What is the
maximum order a word can have? (The order of a word is the power you have to raise
it to in order to get the identity. For example, the order of R is 4.) You can show that
the order of RU, for instance, is 105, by inspecting its cycle structure.

                                       *   *   *

        Where do we go from here? I must mention that I have only scratched the
surface of Cubology in this column. Rubik and ethers are working on generalizations
of various types. There already is a Magic Domino, which is like two-thirds of a
magic cube: two 3X3 layers (see Figure 14-10). You,u can rotate it by q-turns only
about one axis; you must do half turns about the other two. In the START position,
one face is entirely black, the other entirely white, and both faces have the numbers
from 1 through 9 in order. The Domino thus resembles the 15 Puzzle even more
strongly than the cube does.
        Various people have made 2 X 2 X 2 cubes, and such cubes may go on sale
one day. You can make your own by gluing little three-cornered hats over each of the
eight corners of a 3 X 3 X 3 cube. Readers will naturally wonder about such enticing
possibilities as a 4 X 4 X 4 cube. Rest assured-it is being developed in the
Netherlands, and it may be ready soon. Inevitably, there is the question of both higher
and lower dimensionalities. Cube theorists are beginning to discuss the properties of
higher-dimensional cubes.
        The potential of the 3 X 3 X 3 cube is not close to being exhausted. One rich
area of unexplored terrain is that of alternate colorings. This idea was


FIGURE 14-10. Erno Rubik's Magic Domino, scrambled.




Magic Cubology                                                                      327

FIGURE 14-11. Two alternate colorings for the Magic Cube, presenting totally novel
solving problems for the cubist. In both colorings, center orientations do matter.
However, in (a), edge orientations make no difference, and in (b), corner orientations
make no difference.

mentioned to me by various MIT Cube hackers. You can color the cubies in a variety
of ways (see Figure 14-11). Each new coloring presents a different kind of
unscrambling problem. In one variant coloring, edge-cubie orientations take on a vital
importance. In another variant, corner-cubie orientations are irrelevant and centers
matter. Then, moving toward simplicity, you can color two faces the same color,
thereby reducing the number of distinct colors by one. Or you can paint the faces with
just three colors. An extreme would be to have three blue faces meet at one corner and
three white ones meet at the corner diagonally opposite. Inspector General Semah
says that on the cubes he saw, five faces had one color and the sixth face had another
color!
        Who knows where it will all end? As Bernie Greenberg has pointed out:

           Cubism requires the would-be cubist to literally invent a science. Each solver
  must suggest areas of research to himself or herself, design experiments, find
  principles, build theories, reject them, and so forth. It is the only puzzle that requires its
  solver to build a whole science.

      Could Rubik and Ishige have dreamed that their invention would lead to a
model and a metaphor for all that is profound and beautiful in science? It is an
amazing thing, this Magic Cube.




Magic Cubology                                                                                 328

                       On Crossing the Rubicon
                                      July, 1982

                                                          0, cursed spite,
                                                    that ever I was born to set it right!

                                                               (Hamlet, Act I, Scene 5)




THESE days, just "The Cube" will suffice; no one needs to say "Rubik's Cube" to
be understood as making a reference to that great puzzle object. In fact, I have a Cube
in the shape of a sphere, which I sometimes refer to as "the round Cube", but equally
often merely as "that Cube over there". It has been sliced up in the proper way, with
rotating "sides" and an inner mechanism that is the same as Rubik's design. And-what
is even more marvellous -I have. what poses as a Cube but is most definitely not a
Cube: a cubical object sliced in a strange diagonal way, which scrambles in a
devilishly skew manner. Both these puzzles are illustrated in Figure 15-1 The sphere
is, of course, a Cube, while the cube is an impostor in Cube's clothing. (Note: In this
chapter, I use the word "cube" with lowercase `c' as a generic term for any
scrambling-by-rotation puzzle, and with capital `c' to mean the original item: the 3 X
3 X 3 Rubik's Cube.)
        This proliferation of varieties of cube is really an astonishing phenomenon.
Erno Rubik and his somewhat eclipsed Japanese counterpart Terutoshi Ishige began
it, but then it just took off like a prairie fire. Suddenly there were variations on the
Cube turning up all over-little ones, teeny-weeny ones, prettily decorated ones, and so
forth. But in some sense none of these was an essentially different puzzle from the
Cube itself. All of them simply dressed the same internal mechanism in different garb.
The first essentially different cubes I saw came from Japan. They were 2 X 2 X 2's!
One was magnetic, with eight metal cubies sliding around a central magnetic sphere.
The other was plastic, and had an intricate mechanism similar to, but not identical to,
the Rubik-Ishige 3 X 3 X 3 mechanism. It could not be identical, since the keystones
of the 3 X 3 X 3 mechanism are the six face centers-and in, a 2 X 2 X 2, there aren't
any




On Crossing the Rubicon                                                             329

FIGURE 15-1. A cube (a); not a cube (b).

centers! Later I found out that this mechanism is also due to Rubik, and is based on
the 3 X 3 X 3 mechanism. This 2 X 2 X 2, shown in Figure 15-2a, is such a
wonderful, inevitable object-in some ways even more beguiling than the 3 X 3 X 3.
So what puzzles me is: Why aren't they available all over? The 2 X 2 X 2 -Twobik's
Cube?-seems to me an ideal stepping-stone from total novicehood to an intermediate
level of cubistry, as it involves solving only the comers of a 3 X 3 X 3.
        Actually, the 2 X 2 X 2 was not quite the first essentially different cube I
encountered. I had seen a Magic Domino (another Rubik invention-see -Figure 14-10)
much earlier.\_ The Domino is like two of the three layers of a 3 X 3 X 3 cube. Its
square top and bottom layers both can turn 90 degrees, but its four rectangular sides
must turn 180 degrees to allow further moves. Another early variant was the
Octagonal Cube, a cube four of whose edges had been shaved and which, when
twisted, produced some rather grotesque shapes. (See parts (b) and (c) of Figure 15-
2.) Since in this version some of the information about edge parities is lost (you can't
tell whether the "shaved" edges are forwards or backwards in their cubicles), it has
some quirks that make solving it slightly different from solving the full Cube. On .the
full Cube, flipped edges always come in pairs. Here, the same is true except that since
you can't see whether a shaved edge is flipped or not, sometimes you'll wind up with
what appears to be a solved cube, with but a single flipped edge. The first time it can
be quite confusing, if you are used to the full Cube!
        The next variation I encountered was one due to a young German named
Kersten Meier, then a graduate student in operations research at Stanford.




On Crossing the Rubicon                                                             330

He had built a rough working prototype of a Magic Pyramid. It was so rough, in fact,
that it often fell apart as you twisted its sides. Nonetheless, it was clearly an
innovative step, and deserved to be marketed. Hater found out that at nearly the same
time, Ben Halpern, a mathematician at Indiana University, had come up with exactly
the same concept. Both had generalized the Rubik-Ishige 3 X 3 X 3 Cube mechanism
and had seen how e a dodecahedral puzzle on the same principles. Halpern built
ing prototypes of both the pyramid and the dodecahedron. The Meier-Halpern
variations are shown in Figure 15-2, parts (d) and (e).

                                       *   *   *

        As it turns out, Uwe Meffert, another German-born inventor, beat both Meier
and Halpern to the pyramidal punch-but in a different way. Back in 1972, Meffert had
been interested in pyramids and their pleasing qualities when held in the hand.
Somehow, he devised the notion of a pyramid with twisting sides and invented the
concept shown in Figure 15-3. He made a few and found them soothing to play with
and helpful for meditation, but after a while he stored them away and more or less
forgot about them. Then along came Rubik's Cube. Seeing its phenomenal success,
Meffert realized that his old invention might have quite some potential value. So he
quickly patented his design, made arrangements to have his device manufactured in
quantity, and contacted a toy company for the marketing. The end result was the
world success of the Pyraminx, a "pyramidal cube" (in my generic sense of "cube")
that operates completely differently from the Meier= Halpern pyramid.
        Meffert, who now lives in Hong Kong, became deeply involved in the
production and marketing end of his Pyraminx, and began traveling a lot. Through
this he came in contact with other inventors in various parts of the world, and decided
it would be a good idea to market the most interesting toys. of the cube family
worldwide. Among these inventors were Meier and Halpern, and as a result, their
pyramids too will soon be available to puzzle lovers the world over. They will be
known as the Pyraminx Magic Tetrahedron. (I would have preferred "King Tet".) The
dodecahedron will also be available, under the name. Pyraminx Magic Dodecahedron.
(For a catalogue showing Meffert's complete range, write to Uwe Meffert Novelties,
Pricewell (Far East), Ltd., P.O. Box 31008, Causeway Bay, Hong Kong. Incidentally,
Meffert welcomes ideas for new "cubic" puzzles. He also wants to develop a Puzzlers'
Club, in which members would subscribe at a yearly flat rate and receive in return six
or more new puzzles a year. These would be limited editions of particularly complex
or esoteric forms of cubic puzzles. He would like to hear from prospective members.)
        Dr. Ronald Turner-Smith, a friend of Meffert's in the Mathematics Department
at the Chinese University of Hong Kong, has written a charming little book on the
patterns and the mathematics of the Pyraminx,.,




On Crossing the Rubicon                                                            331

On Crossing the Rubicon   332

On Crossing the Rubicon   333

On Crossing the Rubicon   334

called The Amazing Pyraminx, which is available in paperback through Mcffert. In it,
Turner-Smith does for the Pyraminx what David Singmaster did for the Cube in his
Notes on Rubik's 'Magic Cube'. (Incidentally, Singmaster is continuing in his role as
world clearinghouse for Cubology. He now puts out a newsletter amusingly titled
Cubic Circular, available by writing to David Singmaster, Ltd. at 66 Mount View
Road, London N4 4JR, England. Finally, I should mention that a quarterly magazine
called Rubik's will be coming out of Hungary beginning this summer, available for $8
a year. Write to P.O. Box 223, Budapest 1906, Hungary.) Like Singmaster, Turner-
Smith develops a notation and uses it to convey some of the group theory connected
with it, which affords one a deeper appreciation of the object than mere mechanical
solving does.
        It is interesting that there are two distinct ways of manipulating and describing
the action of the Pyraminx. You can rotate either a face or a small pyramid. The two
views are equivalent but complementary, since a face and its opposing small pyramid
make up the whole object. Turner-Smith sees the small pyramids as movable and the
faces as stationary. We shall adopt this view now, and law-.return to comment on the
complementary one. Let us name the four possible moves, then. (See Figure 15-3d.)
Each one rotates a small pyramid, either at the Top (T), Back (B), Left (L), or Right
(R). The letters T, B, L, R stand for clockwise 120-degree turns, and T', B', L', R'
stand for counterclockwise 120-degree turns (as seen when looking at' the rotating tip
along the axis of rotation). Notice that any move leaves all the vertices in place
(although twisted). Therefore, one can consider the four vertices as stationary
reference points, much like the six face centers of the Cube. In fact, at the very start of
the solving process they can quickly be twisted to agree with each other, and from
then on they provide an identifying color for each face. Thus one can consider the
four tip-pyramids either as decorative ornaments or as useful signposts.
        In the Cube, the elementary objects that change location are usually called
cubies or cubelets. What are the corresponding elementary objects here? They are not
all just small pyramids. As on the Cube, it turns out that there are three types: edge
blocks, middle blocks, and the above-mentioned tips. They are shown in Figure 15-4.
As you can see, to each vertex there corresponds one middle block, having three
"trianglets" of different colors, just as does the tip perched on top of it. Also like a tip,
a middle block never leaves its home location, but only twists. As a consequence, the
tips can be considered "trivially solvable" parts of the Pyraminx, and the middle
blocks as "easily solvable".
        This leaves six edge blocks, each having two colors-, that can travel and flip,
just like the edge cubies on a Cube. As a matter of fact, it turns out that the constraints
on flipping and swapping edges are exactly analogous to those applying to the edge
cubies on the Cube: two edges must flip at once, and only even permutations of edge
locations-permutations where an even number of edge swaps have taken place-are
allowed.




On Crossing the Rubicon                                                                 335

FIGURE 15-4. Naming four types of piece in a Pyraminx. In (a), a tip; in (b), an
edge; in (c) a middle block; and in (d), another useful though non-basic unit: a small
pyramid

        This means that one can quickly enumerate the number of different ways
edges can be distributed about the Pyraminx. Without the constraints, the edges could
be dropped into place in 6! (6 factorial), or 720 different ways -the first edge into six
slots, the second into five, and so on. But the requirement that the permutation be
even divides this by two, to give 360. Also, if unconstrained, each edge could be in
either of its two orientations, thus giving 26, or 64, different possibilities-but once
again, we must divide by 2 because of the flipping-constraint, thus getting 32 distinct
flip-states.




On Crossing the Rubicon                                                              336

Multiplying these two figures together, we come up with 11,520 "interestingly
different" states of the Pyraminx. Of course, if you want to take into account the
middle blocks and the tips, each of them has 34 (or 81) ways of twisting, and they are
quite unconstrained, so that you can inflate the figure up to 75,582,720 distinct
scramblings altogether! Perhaps the most realistic figure discounts the tip orientations
but counts the middle blocks. In that case, one has 81 X 11,520=933,120,
"nontrivially distinct" states of the Pyraminx.
        The shortest solving algorithm now known takes 21 twists, and was
discovered with the aid of a computer. It is easy to prove that from some positions one
needs at least twelve twists to get back to START, but the nature of God's algorithm
(which, by definition, always chooses the shortest possible route home) and the
maximum number of twists it requires are unknown, as they are on the Cube.

                                       *   *   *

        When he designed the Pyraminx, Meffert was quite aware that there were.
other ways to slice it up internally, even while keeping the same surface appearance,
with nine trianglets per face. Therefore, he figured out some alternate internal
mechanisms that allow richer modes of twisting. The object 1 have just described is
called the Popular Pyraminx. The Master Pyraminx is a different kind, and is slated to
become available. On it, above and beyond all the movements of the Popular
Pyraminx, each edge can swivel about its midpoint by 180 degrees, thus allowing the
exchange of any two, tips along with the flipping of a single edge piece. (See Figure
15-5.)

FIGURE 15-5. Showing a physically distinct twisting mode, applicable only to the
Master Pyraminx.




On Crossing the Rubicon                                                             337

On Crossing the Rubicon   338

The flexibility requires each middle block: to break up into several pieces as well,
some of which can travel all around the pyramid. Thus one has a much more
complicated puzzle. The mechanism is exceedingly tricky, because during such
swiveling, each of the two moving tips is in contact-with the rest of the-Pyraminx
through a little invisible piece inside the now broken-up middle block. That little
piece does not know to which edge it owes "allegiance". As a result, the invisible
piece and the tip would together fall off (since that contact does not constitute a
permanent link) were it not for a clever piece of engineering that allows each tip to
"lock" its little piece to the appropriate edge piece before the swiveling starts, then to
"unlock" it after the swiveling is over. Turner-Smith cites the number of scrambled
states of the Master Pyraminx as being in excess of 446 trillion.
        Once bitten by the "cube bug", Meffert did not stop here, but moved further
into the world of regular polyhedra. His next step was to design an eight-colored
octahedron each of whose triangular faces is again divided into nine trianglets. How
does it twist? Just as with the Pyraminx, Meffert perceived the possibility of various
modes of twisting. It is interesting that the two equivalent ways of describing the
twists of the Popular Pyraminx .become inequivalent when applied to the octahedron.
Recall that these involved twisting either faces or small pyramids. The reason they
were essentially equivalent is that the rotation of a face is complementary to the
rotation of a small pyramid. However, on an octahedron, rotating a face 120 degrees
is obviously not complementary to spinning a small pyramid (centered on a vertex) 90
degrees. The distinction is shown in Figure 15-6, parts (a) and (b). Realizing this extra
degree of freedom, Meffert designed a mechanism for each of the two ways of
turning.
        The octahedron that will soon be marketed (under the disappointingly clunky
name Pyraminx Magic Octahedron) is the one in which the six small pyramids can
spin. Thus there are three orthogonal axes of rotation just as in the Cube. This
seemingly trivial resemblance to the Cube actually contains much more than a grain
of truth. In fact, the Meffert Octahedron and
the Cube amount to two surface manifestations of one deep abstract idea. To see how
this comes about, notice that a cube and an octahedron are dual to each other: that is,
the face centers of either shape form the vertices of the other shape. Thus the six face
centers of a cube define an octahedron, and the eight face centers of an octahedron
define a cube.
        Imagine a Cube, and, sitting inside it, the octahedron that its face centers
define (see Figure 15-6c). Each twist of a face of the Cube induces a twist on the
corresponding pyramid of the octahedron. Each scrambled position of the Cube seems
thus to correspond to a scrambled position of the Octahedron. But this is not quite
true. To see what is correct, one needs to see what maps onto what, in the
correspondence of Cube and Octahedron. Like the Popular Pyraminx, the octahedron
has tips, middle pieces, and edges. As before, the tips are largely ornamental, and the
middle pieces rotate as wholes. Thus a middle piece on the octahedron (together with
its decorative




On Crossing the Rubicon                                                               339

tip) maps onto a face center on the Cube. This leaves only edge pieces on the
Octahedron-and it is apparent that these, having two facelets, must map onto edge
pieces on the Cube. Where does this leave the Cube's corners? Nowhere. They have
no analogue on the Octahedron, which is a considerable simplification.
        To visualize the Cube-Octahedron correspondence properly, you have to color
one of the puzzles in an alternate manner. Since the Cube is more familiar, let's see
how it has to be altered to "become" a Magic Octahedron. The proper coloring,
corner-centered rather than face-centered, is shown in Figure 15-6d. Stan Isaacs, a
computer scientist and puzzlist par excellence, has made up one of his dozens of
cubes to simulate a Meffert Octahedron. Someone fluent in solving the ordinarily
colored 3 X 3 X 3 Cube will therefore find that their expertise does not quite suffice
to handle Isaacs' strangely colored cube, because now the orientation of face centers
matters! On the other hand, there is a corresponding simplification as well: "quarks"
no longer exist on this cube. That is, there is no such thing as a twisted corner, simply
because all the corner cubelets are white on all sides.
        All you need to solve this cube (or the Octahedron) is the ability to restore the
edges and face centers (with the added novelty of orientations). Of course, not all
"magic octahedra" will be equivalent to simple recolorings of the 3 X 3 X 3 cube,
since they may not turn about those three axes. In particular, Mcffert's alternate
twisting-mode for the octahedron (where faces twist 90 degrees) is quite unrelated to
the Cube.
        In his 1982 catalogue, Meffert shows a picture of an icosahedron (guess what
its name is!) whose twenty triangular faces are not subdivided at all; they move five at
a time, swirling about any of the twelve vertices. (See Figure 15-2f) Since the
movement is vertex-centered rather than facecentered, it should make you think of the
icosahedron's dual solid, the dodecahedron. The dual puzzle would have face-centered
movement, in the same way as the dual puzzle to the Octahedron, with its vertex-
centered movement, is the Cube, with its face-centered movement. (Incidentally, what
would be the dual puzzle to the Pyraminx?}
        In fact, in Meffert's catalogue are shown two other dodecahedral puzzles,
reproduced in Figure 15-2g and Figure 15-2h, for your amazement and bemusement.
The less complicated one with the asymmetric-looking slices is called the Pyraminx
Ball, and the beautifully crisscrossed one is called the Pyraminx Crystal. The Ball has
four axes of rotation, like the Pyraminx, while the Crystal has six. These should be
hitting the market in midsummer.

                                       *   *   *

         At this point, you might well be wondering whether there could be a cube -I
mean a genuine, six-sided, square-faced cube!-with a vertex-centered twisting
mechanism. No sooner said than done! Tony Durham, a British journalist, was the
first to think of this idea. He showed his design to Mêffert




On Crossing the Rubicon                                                              340

FIGURE 15-7. Tony Durham's Skewb, caught in mid-twist (a). In (b), the labeling of
the' Skewb''s eight corners. (See also Figure 15-1.)

who developed it into a marketable product by incorporating mechanical features that
had proved useful on the Pyraminx. The object in question is shown at rest in Figure
15-lb and in motion in Figure 15-7a. I call this the Shewb, although Meffert gives it
the more prosaic title of Pyraminx Cube.
        Each of the Skewb's four cuts slices the whole into two equal halves. Each cut
perpendicularly bisects one of, the four spatial diagonals of the cube. If you think
about it, you will see that the shape traced out by each cut as you run around the
cube's surface is a perfect hexagon. Each cut crosses all six faces, so that every turn
affects all the faces at once. In this respect, the Skewb is more vicious than the Cube,
where on each turn two faces are exempt from change. Despite the simplicity of this
object, it is quite hard to get used to its skew twist. Of course, that is part of its charm.
        Durham offers some insightful commentary on his invention in a remarkable
set of notes he has written entitled "Four-Axis Puzzles". I would like to quote a few
paragraphs from this document.

           The symmetry group, generated by four threefold axes is the rotation
   group of the tetrahedron, and has order twelve. Almost all the well-known
   polyhedra, regular as well as semiregular, possess this tetrahedral symmetry,
   though their own symmetry may be much richer. So a four-axis mechanism may
   be put inside a polyhedral puzzle of any regular or semiregular shape, and the
   puzzle will keep its shape during play.. The Pyraminx Ball may look odd at first
   glance, but it illustrates the beautiful way in which tetrahedral symmetry is
   buried in the richer symmetry of the dodecahedron.
           The cube mechanism found by Rubik does not have this property. It uses
   fourfold rotation axes, which are generally found only in the cube/octahedron
   family of solids. Thus, it is possible to 'build out' a Rubik cube into the shape of
   a dodecahedron. But to preserve that shape during play you must restrict




On Crossing the Rubicon                                                                 341

  yourself to half-turns. Quarter-turns invoke a symmetry which the
  dodecahedron does not possess.
          All four-axis puzzles have a central ball or spindle. Four pieces (usually
  corners) are pinned directly to the ball. The standard ' Pyraminx has six free-
  floating edge pieces with `wings' that hook under the corner pieces. The
  analogous free-floating pieces on the Pyraminx Cube are the square face-
  centers. The four-faceted pieces on the dodecahedral Pyraminx Ball play the
  same role.
          The Pyraminx Cube and Ball have four more free-floating pieces, which
  again are corners. These pieces have their own 'wings' which, in the START
  position, hook under the first set of free-floating pieces. Thus, there is a three-
  level hierarchy of interlocking pieces, conceptually similar to Rubik's, but
  geometrically very different.
          All eight corners of the Pyraminx Cube look alike. At first sight one
  might think that any two corners could be made to change places. In fact, four
  of the corners are free-floating and four are rigidly fixed to the central ball. The
  two types can never change place. The square shape of the face center pieces is
  deceptive, too. Inside, the mechanical parts of the square pieces are not so
  symmetrical. Such a piece can never return to its starting position (relative to
  the rigid set of four corners) rotated by 90 degrees. Only half-turns are possible.
          The standard Pyraminx has obvious fixed points-the four corners.
  Confronted with a Pyraminx Cube and knowing that four corners are fixed and
  four are free, one naturally wonders which are which. Actually it makes no
  difference. The four free corners move independently of the fixed ones, but they
  always move together as if physically linked.

        Durham proceeds to give Turner-Smith's TBLR-T'B'L'R' notation for the
Pyraminx, and mentions that it is adaptable to any four-axis puzzle (such as his
Skewb), simply by letting TBLR name four of the centers of rotation. (On the
Pyraminx, this could mean either the four tips or the four face centers. On the Skewb,
this would be four of the tips, leaving four other tips unnamed. See Figure 15-7b.)
Then any move can be transcribed. If it is centered on one of the named spots, just use
the proper notation. If it is centered on one of the four unnamed spots, use the name
for the complementary move, since it doesn't matter which half of the puzzle twists.
(You may want to think about that for a moment. Actually, it is obvious, but it sounds
like a tricky point.) Durham points out that it is sometimes useful to have names for
the four remaining spots and for twists around them. He lets t, b, 1, r fulfill that
purpose. Thus T and t accomplish the same thing internally to the puzzle, but they
leave it hovering in space in a different overall orientation. Although he concedes that
it may become confusing, Durham advocates using a mixed notation on occasion.

          Sometimes you need to mix the notations to see what is going on. TbT'b'
  is one of the useful class of moves called commutators (two moves followed by
  their inverses-thus of the form xyxy'), though you would never guess so from its
  description in regular coordinates (TBL'B') or alternate coordinates (tlt'b').




On Crossing the Rubicon                                                              342

           The Pyraminx Cube and Ball may be described as deep-cut puzzles in
   contrast to shallow-cut puzzles such as Rubik's Cube.. In the latter, the cuts are
   made near to the surface. In deep-cut puzzles, they slash close to the puzzle's
   heart. The bulk of a shallow-cut puzzle remains stationary while you turn a
   small part of it. A deep-cut puzzle, however, raises serious doubt as to which
   part has been turned and which has remained stationary. This is why alternate
   sets of coordinates have to be taken seriously on deep-cut puzzles.
           Deep-cut puzzles also dictate a 'global' approach to solution. It is
   peculiarly difficult to work on one area of the puzzle without affecting the rest.
   However, as solution proceeds, this very fact comes to your aid. Pairs of corners
   magically untwist in synchrony. The last flip, the last swap is done for you
   automatically. As you close in for the kill, billions of pathways down which the
   puzzle might escape are closed off to it. Parity constraints are at work, and when
   every move activates five or eight interlocked permutation cycles-as it does in a
   deep-cut puzzle-parity constraints are powerful.

       In the section of his notes having to do with parity constraints, Durham.
includes the following humorous but insightful apology:

   Please forgive the loose use of the term parity to include tests for divisibility by
   3 (not only 2) or even more distant concepts. We shall We the term parity
   restriction for any constraint on imaginable transformations of the puzzle that
   prevents their accomplishment in normal operation of the puzzle. The list does
   not, for example, include the rule: 'Thou shalt not swap a face piece with a
   corner piece.' It is just too far-fetched. One might as well try to imagine a move
   that transformed the entire puzzle into depleted uranium or Gorgonzola cheese.

Then he lists all the Skewb's "parity" constraints, in his generalized sense of the term.

   1. The four (fixed),corners TBLR may be permuted among themselves, as may
      the remaining four corners tblr, but mixing between the two sets is
      prohibited.
   2. TBLR themselves move as a rigid tetrahedral unit. This constraint applies to
      their positions in space only (not to their orientations).
   2a. For exactly the same reasons, the remaining four (free) corners tblr move as
      a tetrahedral unit. They move independently of TBLR. In fact any of the
      twelve possible relative positions of tblr and TBLR can be reached in at most
      two puzzle moves.
      Although TBLR are fixed and tbir are free-floating, mathematically
      speaking, 2 and 2a have exactly the same status. Writers on the Rubik Cube
      have generally regarded the transposition of two face centers as an
      'unimaginable' transformation, while the swapping of two edge pieces is
      'prohibited but imaginable'. By analogy with this convention, 2a counts as a
      parity restriction while 2 does not! This is plainly unsatisfactory, and a better
      and more precise definition of 'parity' is badly needed. Is it a question of
      geometry? Of mechanics? Of topology? Note that the problem is in
      enumerating the impossible positions. The possible positions are readily -
      counted.



On Crossing the Rubicon                                                               343

  3. The sum of the twists of corners TBLR is always equal, modulo 3, to the
      twistedness of the puzzle, taken as a whole.
  (Here, twist applies to corners, and is either 0, + 1, or -1. A corner's twist is
      measured relative to the rigid tetrahedron to which it belongs. Thus the twist
      of T is measured relative to TBLR. A clockwise rotation of a corner counts
      as + I, counterclockwise as - 1. By contrast, the twistedness of the puzzle as a
      whole is a function only of the positions of the corners, not of their
      orientations. If the relative positions of TBLR and tblr are as in the START
      position, then the twistedness is 0. If they can be restored to START by one
      clockwise puzzle move, the twistedness is -1, and if by one counterclockwise
      move, then + 1. If it takes one of each type, then the twistedness is again 0.)
  3a. Same as 3, only with tblr.
  From 3 and 3a, it follows that the total twist of TBLR always equals the total
      twist of tblr. Also, it follows that it is impossible to turn a single corner by
      120 degrees (i.e., to create an isolated quark). One might paraphrase 3 and 3a
      by saying that the puzzle 'knows', in three distinct ways, how many turns it is
      away from START (modulo 3).
  4. It is impossible to transpose exactly two face pieces.
  5. It is impossible for any face piece to turn in place by 90 degrees. 6. It is
      impossible to flip a single face piece through 180 degrees.

Durham offers proofs of these interesting facts, but as they are for the most part
analogous to those on the Cube, I shall omit them here. Bycombining all these
constraints, Durham comes up with the total number of '-scrambled states of his
Skewb, which is 100,776,960. However, this assumes you have a way of telling the
orientation of a face center, which (unless you mark. it up) you don't. Hence the
number of visually distinguishable states is reduced by five factors of two, to
3,149,280-a rather smaller number than for the Cube (4 X 1019), but certainly the
difficulty does not scale down proportionately with the number of states. (Could
you even imagine what it would mean for a puzzle to be "ten trillion times easier"
than Rubik's Cube?)

                                     *   *   *

Durham's final observations carry Solomon Golomb's beautiful analogy between
cubological phenomena and those of particle physics to even greater heights.
Golomb pointed out that many fundamental particles have their counterparts on the
3 X 3 X 3 Cube. They include the quarks (q), antiquarks (q), mesons (qq pairs),
baryons and antibaryons (qqq and qqq trios). Durham extends the analogy as
follows:

     The definition of twist must be modified for the purpose of particle physics.
     A clockwise twist of one of the corners TBLR is now given the value + 1/3,
     as is a counterclockwise twist of any of the corners tblr. Either of these is a
     quark. Its opposite is an antiquark with value -1/3. It will be seen that twist
     corresponds to baryon number. The total twist of all corners is always an
     integer. A single puzzle move is always a meson




On Crossing the Rubicon                                                              344

     Quarks at the corners TBLR will be regarded as 'up' or u quarks; those at tblr
     will be 'down' or d quarks. Both quarks have isotopic spin 1/2. They are
     distinguished by the orientation of the isospin vector in its abstract space.
     The projection of the isospin, I, , has the value + 1/2 for the u quark and -1/2
     for the d quark. In the absence of strangeness, charm, etc., the electric charge
     Q of a particle is given by Q=Ix +B/2, where B is the baryon number. So u
     quarks have charge 2/3, while d quarks have charge - 1/3. (All the quantum
     numbers are multiplied by -1 for the antiquarks.) Again the puzzle models an
     important feature of observed reality: all particles have integral electric
     charge.
     The relevant quantum numbers for our two quarks are as follows:

                              U      D
                      B      1/3     1/3
                      1      1/2     1/2
                      Ix     1/2     -1/2
                      Q      2/3     -1/3

     We can now assemble various hadrons (strongly interacting particles), as
     shown in the table below. Each particle is represented by two rows having
     four symbols each. The four places in the top row represent the twists on the
     TBLR corners; in the bottom row the same is done for the tblr corners. A
     quark is denoted by '+', an antiquark by '-'.




       Isotopic' symmetry is a global symmetry, and the strong (nuclear) force is
     invariant under transformations that rotate the isotopic spin vector by the
     same amount for all particles. Such a transformation would, in a continuous
     fashion, transform all u quarks into d quarks and vice versa. Protons and
     neutrons would swap roles. The analogous process for the puzzle is the
     continuous rotation of the whole puzzle in space. It can indeed bring the
     TBLR corners to the former position of the' tblr corners, so that an up quark
     becomes a down quark.
       This makes no difference to the 'strong interaction' (i.e., the normal
     operation of the puzzle). The TBLR and Iblr corners are functionally
     identical. But it. matters, if you try to dismantle the puzzle: you will find that
     one set of corners is fixed to the core, and one is not. Such dismantling
     operations can be thought of as weak or electromagnetic interactions, which
     can break the conservation rules obeyed by the strong interaction. Actually
     they break the rules rather too well, since they allow the creation of single
     free quarks.

     Durham points out that the analogy still has weaknesses, such as the facts
     that neither charge nor baryon number is conserved, that there is no

On Crossing the Rubicon                                                               345

analogue to spin, that only two "flavors" of quark are represented (up and down),
and that quark "color" is not modeled. Golomb, in the meantime, has been actively
trying to find a way of modeling quark color in the 3 X 3 X 3 Cube analogy.
Whatever the failings of this analogy, I find it one of the most provocative .of all
analogies I have ever encountered anywhere, and will be most astonished if it is
purely coincidental. I somehow cannot help but believe that the fascinating
patterns shared by these macroscopic puzzles and the microscopic particles reveals
some underlying order and set of principles common to both. Indeed, I have faith
that, if looked at in the proper way, the group-theoretical principles that govern
these parity constraints on "cubes" can be transferred to the domain of particle
physics, and yield fresh insights about the reasons for the symmetries among
particles. There! If that doesn't prod some particle physicist into looking into this, I
don't know what will!

                                          *   *   *

Perhaps my favorite "cube" is the one I dubbed the IncrediBall. It is due to a
German educator from Dortmund named Wolfgang Kuppers, and is in Meffert's
catalogue. As of the time of this writing, I may be the world's fastest IncrediBall
solver (or at least the fastest on my block!), with an average time of about six
minutes. However, I am sure that my glory will not last long, once this puzzle is
marketed widely by the Milton Bradley Company sometime this summer. Their
trade name for it will be ,l assi*Ball. It is pictured in Figure 15-8.
       This I-Ball is basically a rounded-off dodecahedron each of whose twelve
faces (dodecalets, I'll call them) has been subdivided into five elementary
"trianglets". Thus there are 60 such trianglets. If, instead of seeing them in groups
of five, you take them three at a time, you'll find that they define a rounded-off
icosahedron (the dual of the dodecahedron). Such a group of three trianglets I call
an icosalet, and there are twenty such, each one having a unique arrangement of
three colors. The icosalets are the elementary, unbreakable units out of which the
IncrediBall is constructed; they correspond to the cubelets on the Cube, or the
elementary pyramids of the Pyraminx. Whereas on the Cube there are three kinds
of cubelet (edges, faces, and corners), here all icosalets are of a single type. For
this reason, the I-Ball is less forbidding than at first it might appear. Its pristine
state is one in which each dodecalet is all of one color. Meffert has used only six
colors, rather than twelve, each color being used in two antipodal dodecalets, but
this does not in any way change the difficulty of the puzzle.
       The way it turns is a little surprising. Any group of five icosalets that meet at
a point (the center of a dodecalet) form what I call a circle, which will rotate as a
unit, twisting 72 degrees to the left or right. (Such a circle is analogous to a
"layer"-a face together with its fringe-on the Cube.) Thus five such




On Crossing the Rubicon                                                                346

FIGURE 15-8. Wolfgang Kuppers' IncrediBall (or Impossi*Ball, if you wish). In
(a), the pristine state. The triangles with curved sides are called icosalets. In (b),
an IncrediBall caught in the midst of a "bumpy twist" - Each such twist involves
rotating a "circle" (composed of five,,: icosalets) through 72 degrees. In (c), a state
with just one quark visible (one icosalet twiste ~120 degrees clockwise). In (d), one
icosalet has been removed. This allows another icosalet to s in and occupy the
vacuum, meanwhile leaving behind its own vacuum. As an icosalet-shaped glides
around the puzzle, order can be created or destroyed This sphere-based puzzle
thus c resembles Sam Loyd's planar "15 puzzle'




On Crossing the Rubicon                                                               347

twists return that group to its starting position. However, the "circle" defined by the
five icosalets is not truly circular, and if the trianglets were rigidly held at a fixed
distance from the center, it simply would not be possible to rotate such a group.
But Meffert's mechanism ingeniously gets around that problem by having the
icosalets lift up slightly as they go over 'bumps", so that the solid flexes noticeably.
As a result, twisting the I-Ball has a delightful "organic" feel to it.
         The constraints here are the same old story: all permutations are even,
which means you cannot swap two icosalets-the best you can do is cycle three of
them, or swap two pairs simultaneously; and of course, quarks and antiquarks must
add up to a total twist that is integral. Taking into account these constraints, I
calculate that the total number of IncrediBall scramblings is
23,563,902,142,421,896,679,424,000, or 24X 1024-about 24 trillion trillion. This
is not quite a million times larger than the figure for the Cube. It's also about 40
times larger than Avogadro's number, for whatever that's worth.
         How hard is it to solve this puzzle? Is it harder than the Cube? I found + t
easier, but that's hardly fair, since I had already done the Cube. However, in
Durham's terms, the IncrediBall is decidedly a "shallow-cut" puzzle, which means
that a more or less local approach will work. I found that, when I loosened my
conceptual grip on the exact qualities of my hard-won operators for the Cube, and
took them more metaphorically, I could'transfer some of my expertise over from
Cube to I-Ball. Not everything transferred, needless to say. What' pleased me most
was when I discovered that my quarkscrew" and "antiquarkscrew" were directly
exportable. Of course, it took a while to discover what such an export would
consist in. What is the essence of a move? Which aspects of it are provincial and
sheddable? How can one learn to tell easily? These are very difficult questions, to
which I do not have the answers.
         I gradually learned my way around the IncrediBall by realizing that a
powerful class of moves consists of turning only two overlapping "circles" in a
commutator pattern (xyx'y'). So I studied such two-circle commutators on paper, as
shown in Figure 15-9, until I found ones that filled all my objectives. They
included quarkscrews, swaps, and 3-cycles, which form the basis of a complete
solution. In doing this, I came up with just barely enough notation to cover my
needs, but I did not develop a complete notation for the IncrediBall. This, it seems
to me, would be very useful: a standard 'universal notation, psychologically as well
as mathematically satisfying, for all cubelikeles. However, it is a very ambitious
project, given that you would have to anticipate all conceivable future variations on
this fertile theme-hardly a trivial undertaking!
It is interesting that my diagrams of overlapping circles turn out to be closely
connected with another lovely family of generalizations of the Cube, due to a
Spanish physicist named Gabriel Lorente. His puzzles are mostly planar and
consist precisely in networks of overlapping circles. (See Figure




On Crossing the Rubicon                                                                348

FIGURE 15-9. Operators involving just two overlapping circles. In (a), names for
the four possible 72-degree twists. In (b), the operator "doui " (down-out-up-in) is
applied. Note that it has the form of a commutator, involving alternating inverses.
In (c), the outcome is summarized: a double swap has been effected.

15-10.) The planar ones he calls the Grill and the Trebol. In each of them, circles
can be given partial twists and pieces of them are thereby shuffled and
redistributed. Extending this notion to a spherical surface, Lorente came up with an
elegant IncrediBall-like puzzle, which he calls the Florid Sphere.
         When you look closely at Lorente's puzzles, the IncrediBall, and even the
Cube, you begin to see that the essence of all these puzzles seems to reside in
overlapping orbits. In fact, one could even maintain that the three-dimensionality
of all these puzzles is irrelevant; their interest is essentially due only to the
properties of intricately overlapping closed orbits in a two-dimensional space,
possibly curved like a sphere.
         The quintessential planar overlapping-circle puzzle was invented, as it turns
out, way back in the 1890's, although recently it has been repeatedly rediscovered
in the wake of the Cube. All such puzzles basically involve two circles of marbles
that intersect at various spots. (See Figure 15-11.) You can choose to cycle either
circle, and the marbles at the intersections will thus be absorbed. into whichever
circle is moving.
         While we're discussing two-dimensionality, it is worthwhile pointing out




On Crossing the Rubicon                                                              349

FIGURE 15-10. Four puzzles by Gabriel Lorente. In (a) and (b), two schemes he
calls "Grills " Note that both are based on a square lattice of circle centers. In (c),
his "Trebol "puzzle, where centers form a triangular lattice. In (d), the centers of
circles lie on a sphere. This is his "Florid Sphere" : Which previously discussed
puzzle is it equivalent to?

that the Incredi\&all's internal construction allows it to be transformed rather
amazingly into what I call the "19" puzzle-a two-dimensional curved-space version
of Sam Loyd's famous "15" puzzle (the 4X4 square puzzle with one "squarelet"
removed, allowing you to rearrange the remaining 15 squarelets by shifting the
hole about). This was first observed by Ben Halpern, while he was idly playing
with an IncrediBall. He had removed one single icosalet (which is possible, one of
the beauties of the IncrediBall being that its mechanism readily allows disassembly
and reassembly),, leaving a hole, and he observed that, because all icosalets are




On Crossing the Rubicon                                                               350

FIGURE 15-11. Does this 90-year-old type of two-dimensional puzzle, with just
two intersecting rings of marbles, capture the ultimate essence of all modern
"cubic" puzzles?

congruent, the hole could wander about all over the sphere, just like the square
hole in the 15 puzzle. (See Figure 15-8d.) Again, this seems to underscore the two-
dimensional nature of these puzzles.
        The claim that these puzzles are two-dimensional comes from the fact that
only pieces on their surfaces move; there is. no exchange between the interior and
the exterior. For an extreme case, imagine the Earth as a giant puzzle, its entire
surface covered with trillions of overlapping circles of marbles. With a hundred
million turns, you could ship a marble from New York to San Francisco. Clearly
this would be in essence a two-dimensional puzzle. The smallness of the circles
relative to the size of the Earth makes this obvious. (However, I surely wouldn't
want to think about solving such a puzzle, whether it's two-dimensional or not!)
        By contrast, consider two objects about to come out: Ideal's 4 X 4 X 4 cube,
tastelessly marketed as Rubik's Revenge, and Meffert's Pyraminx Ultimate, a 5 x 5
X 5 with shaved corners. Both are shown in Figure 15-2, parts (i) and, (j). In these
objects, there are circles on a much more global scale. Namely, the 4 X4 X4 has an
"Arctic Circle",, a "Tropic of Cancer", a "Tropic of Capricorn", and an "Antarctic
Circle". The Pyraminx Ultimate has an Equator as well.




On Crossing the Rubicon                                                            351

         O the S X 3 X S Cube, one could get affray with' ignoring the Equator by
describing equatorial twists in terms of their complements, like rotating the slices
of bread instead of the meat in a sandwich. Singmaster's notational choice for the 3
X 3 X 3 Cube reflects his propensity to describe face centers as stationary. Thus
for him, bread slices move while the meat stays put. Theoretically, this is fine, but
realistically, people just do not hold their sandwiches-pardon me, their Cubes-in
one fixed orientation. Moreover, when you pass to higher orders, this view will not
suffice. Imagine a multilayer club sandwich with three slices of bread and two
different kinds of meat. For this, you simply have to expand your notational
horizons!
         An elegant set of names for the six possible meat-slice, or equatorial,
moves on the 3 X 3 X 3 Cube has been suggested by John Conway, David
Berlekamp, and Richard Guy in their book Winning Ways. They employ Greek
letters with clever mnemonic justifications. These are shown in Figure 15-12. With
some modification, they could be adapted to slices on higher-order cubes.
         Slice moves of this more global sort are like giant circles of marbles
stretching around the Equator or the Tropic of Capricorn; their radii are of the
same order of magnitude as the radius of the underlying threedimensional object.
The topology of linkage of circles becomes much more complicated than in the
case where the circles are"small and every connection is very local. To describe the
linkage economically, one would be forced to talk about the way the circles are
embedded in 3-space. In this sense, the higher-order cubes can truly be said to be
intrinsically three-dimensional puzzles.
         There are, it seems, endless new spinoffs of the Cube being created. It is




       FIGURE 15-12. The Conway-Berlekamp-Guy nomenclature for twists of a
3 X 3 X 3 Cube's equatorial slices. This notation can be generalized to higher-
order cubes.




On Crossing the Rubicon                                                             352

a very fertile idea. H. J. Kamack and T. R. Keane, both chemical engineers, sent
me a beautiful paper in which they describe their simulation of a four-dimensional
3 X 3 X 3 X 3 cube on a computer. They call it Rubik's Tesseract, and they have
computed the number of possible states it has. That number is: the product of 24!
X 32! X 16!/4 (the number of permutations of position of the elementary "tessies
out of.which it is built) with (224/2) X (632/2) X (1216/3) (the number of legal
orientations of the tessies within their niches, which the authors somewhat
hesitantly term "tessicles", by analogy with "cubicles"). This number comes to
approximately 1.76 X 10120, which, they point out, is about the same size as the
number of possible games of chess. (I don't think it would be an exaggeration to
say that if Ideal were marketing this puzzle, their publicity would sharr1lessly
proclaim, "Over 3 trillion combinations!") Kamack and Keane have made many
provocative discoveries, which unfortunately I have no space to report on at this
time. I was also sent a fascinating paper by George Marx and Eva Gajzago, two
physicists at the famous Roland Eotvos University in Budapest. In it they give a
definition of "entropy" on the Cube and. describe some statistical results computed
by a grammar school student named Victor Zambo. These are matters I would like
to go into at some future time.

                                      *   *   *

        I would like to close by discussing the astonishing popularity of the Cube.
In the New York Times' paperback bestseller list of November 15, 1981, three
cube booklets figured on the list. The positions they occupied? First, second, and
fifth. People often ask, "Why is the Cube so popular? Will it last? Or is it just some
sort of fad?" My personal opinion is that it will last. I think that the Cube has some
sort of basic, instinctive, "primordial" appeal. Its conceptual pizzazz comes from
the fact that it fits into a niche in our minds that connects to many, many general
notions about the world. So here is an. attempt to characterize that quality.

        • To begin with, the Cube is small and colorful. It fits snugly in the hand
and has a pleasing feel. Twisting is a fundamental and intriguing motion that the
hand performs naturally. The object itself has overall symmetry, so that it can be
rotated as a whole without its "feel" changing. (This is in contrast to many puzzles
that have at most one axis of symmetry.) Quite surprisingly, there are not many
puzzles or toys that give the mind and fingers a genuine three-dimensional
workout.
        • Although it gets all scrambled up, the object itself stays whole. (This is in
contrast to many Humpty-Dumpty-ish puzzles that come apart into. scads of pieces
that may get scattered all over the floor.) That it . manages to stay in one piece
when it has so many independent ways of twisting is initially amazing, and
remains mysterious even after you've seen its "guts".




On Crossing the Rubicon                                                               353

        • T'he object is a miniature incarnation of that subtle blend of order and
chaos that our world is. Most of the time, you just cannot predict what
repercussions even simple actions will , have-they simply have too many side
effects. A few tiny actions can have vast, interlocking consequences, and become
practically un-undoable. One can easily become paralyzed by fear, not wanting to
make any move at all, sensing that with no trouble at all one can get totally,
irretrievably, hopelessly lost.
        • There are plenty of patterns, some attainable, some unattainable.
Sometimes they are simple to generate, but one' can't see how they emerge.
Sometimes they are hard to generate, yet one clearly understands how they arise.
        • There are many routes to any state, and the shortest is nearly always
completely unknowable. The solution to a difficult situation is hardly ever to back
out the way you came in, but to find an alternate and completely different escape
route. One feels a little like someone trapped in a cave with no light, unable to
sense the whole space, able to grope about only very locally, wondering whether it
is even humanly possible to have such an overview. (One wonders about God's
algorithm:\_ Is it humanly comprehensible?)
        • The Cube is a rich source of metaphors. It, furnishes analogies to particle
physics (quarks, etc.), to biology (a move-sequence as a "genotype" and the pattern
it codes for as a "phenotype"), to problem solving in everyday life (breaking a
problem into parts,. solving it stage by stage), to entropy and path-finding, and on
and on. It even touches theology ("God's algorithm") and many other phenomena.
        • There are different approaches to understanding the Cube. In particular,
there is a strong contrast between the "algebraic" approach and the "geometric"
approach. In the algebraic, or mathematical, approach, long sequences of
operations are compounded out of shorter sequences,, so that after a while one has
no idea why one is doing the various individual twists-one just relies on the
sequences, as wholes, to work. Though efficient, this is risky. In the geometric, or
commonsense, approach, eye and mind combine to choose twist after twist, each
twist having a clear reason as part of a carefully charted pathway. Though
inefficient, this is reliable. These approaches, of course, can serve as metaphors for
styles of attacking problems in life.
        • The Cube's universe has a strange population. Aside from its varieties of
"cubies and modes of twisting, there are such intangible qualities as "flippedness"
or "twistedness", which one quite literally moves about on the Cube (e.g., in the
form of quarks), just as one moves the tangible cubies Similarly, the word "here"
can designate a "place" that moves to and fro during a sequence of twists. The
interlocking and nested reference frames that one jumps between in trying to
restore order to the Cube vividly exemplify the layered way in which we




On Crossing the Rubicon                                                              354

conceive of space-indeed, `the layered 'may in which concepts themselves are
structured in-our minds.
        * Among the Cube's less intellectual charms are the magic of motion too
swift for the eye, the thrills of speed, competition,, and grace; the varying levels of
knowledge that one can gain, the enjoyment of exchanging information and
insight. And, needless to say, the very idea that such a tiny innocent object
conceals such a vast universe\_ of potential.
        * Finally, consider the metaphor the Cube offers for the state of the world
(one that has been exploited in various political cartoons). The globe is in a mess
(as shown in Figure 15-13), and the leaders of various




FIGURE 15-13. The sad state of the globe.

countries want it to be "fixed". But they are unwilling to relinquish any tiny bit of
order they have achieved. They cling to old, useless achievements because they are
too fearful of letting go and temporarily abandoning what partial order they have in
order to achieve greater order and harmony. They lack a mature, global view, one
that recognizes that a willingness to make sacrifices in the short run can wind up
producing much greater gains in the long run.

       I am confident that The Cube, as well as "cubes" in general, will flourish. I
expect new varieties to appear for a long time to come, and to enrich our lives in
many ways. It is gratifying that a toy that so challenges the mind has found such
worldwide success. I hear that it's now very popular in China.




On Crossing the Rubicon                                                               355

Perhaps one day it will even penetrate into the Soviet Union, to - my knowlege the
last bastion of the Cube-Free World.

Post Scriptum
        I wrote the preceding two columns over a year apart. It has now been two
years since the second of them was written. The major cube news since then is, sad
to say, that there has not been much major cube news since then. What apparently
happened was simply a worldwide cube glut. Cubescubical and otherwise-were
coming out of everybody's ears, and it was just a little too much. I can understand
that, but it saddens me to see something so exciting fade so totally.
        There are still a number of things worth mentioning. A good place to begin
is with the origin, of the cube-that is, of magic solids in general. Shortly after my
second cube column appeared, I received a rather plaintive letter from a Fresno
high school teacher named William 0. Gustafson, who claimed that he was, in
some sense, the true inventor of the idea of the Magic Cube. What he actually had
invented-in 1958-was a sphere sliced by three orthogonal planes into eight
congruent pieces (octants), in such a way that any two opposite hemispheres (each
composed of four octants) could turn. This amounts to a spherical 2 X 2 X 2 cube-
a- cubical variant of which was marketed by Ideal Toy Company some twenty-odd
years later under the name "Rubik's Pocket Cube". Gustafson called his toy
"Gustafson's Globe".
        To substantiate his claim, Gustafson enclosed photocopies of a good deal of
correspondence he conducted in 1960 with numerous toy companies (he wrote 76
of them!), the Japanese patent office (he received a patent)-and even Martin
Gardner. Gardner's card to him was interesting. It said:

  That is an interesting puzzle that you propose, but I am at a loss for suggestions
  on how to interest a toy dealer in it. My experience has been that it is almost
  impossible to make any money with a puzzle unless you are the manufacturer
  yourself, with your own toy company.

An interesting comment, in light of what happened with Rubik's Cube. In addition
to his 2 X 2 X 2 Globe, Gustafson developed a 3 X 3 X 3 version, but - felt that he
should work first on getting the simpler puzzle out, so most of his correspondence
concerned that one.
        Gustafson also enclosed for me a photocopy of a wry letter of condolence
sent to him by a former student who, when he encountered Rubik's Cube, vividly
recalled Gustafson's Globe from decades earlier. The card read: "With sincere
sympathy in your recent loss, and a hope that time has helped in some small way to
ease the sorrow in your heart". Below those poetic




On Crossing the Rubicon                                                             356

lines were the words "Gustafson's Globe", crossed out, and then the words "Rubik's
Cube".
       I did a little bit of checking around, including talking with David
Singmaster, probably the world's leading Cubological and Cubohistorical
authority, and discovered that there is something to Gustafson's claim of priority.
Not that it is likely that ErnO Rubik or Terutoshi Ishige ever heard of Gustafson.
Nor is it by any means certain that Gustafson's Globe, had it been picked up by
some toy company, would have been the overnight sensation that the Rubik-Ishige
Magic Cube was. Still, though, it seems only fair to point out that people besides
Rubik and Ishige had smelled some of the same alluring aromas in previous years,
and for various reasons had not been able to arouse the interest of the world.

                                      *   *   *

         It is one of my firmest beliefs that good ideas almost never come out of
nowhere, and that if a good idea arises in one person's mind, it is almost sure to
have arisen in someone else's mind in some closely related version, or to do so
very shortly. For that reason, whenever I am writing about a discovery or
invention, I always try hard to indicate multiple credit when I can discover the
people to whom genuine credit is due. The trouble is, when you bend over
backwards to be equitable (notice how I always talk about Ishige and Rubik in the
same breath, for instance), what inevitably happens is that someone you
inadvertently slighted then writes you with some mixture of indignance,
consternation, and disappointment, and requests equal time.
         I am glad to mention Gustafson's name and to give him credit for having
had perhaps the world's first insight into this kind of three-dimensional rotational
puzzle. But at the same time, I do not wish to leave out the names of Frank Fox, a
British inventor who in about 1970 discovered-and patented-a 3 X 3 X 3 twisting
sphere, and Larry Nichols of Cambridge, Massachusetts, who invented and
patented a 2 X 2 X 2 cube around 1972, and who has just won a suit against Ideal
Toy Co. for not giving him royalties on his invention. Whether Nichols' claim is
any more deserving of retroactive compensation than those of Fox or Gustafson, I
am not competent to say.
         All I can say is, these things get very, very messy-particularly when large
amounts of money (or glory) are concerned. In the case of all three inventors-
Gustafson, Fox, and Nichols-it seems clear that their inventions were far flimsier
than the Rubik-Ishige Cube, and that the real reason the Rubik-Ishige Cube took
off was that it could be manufactured and that it did hold together. But perhaps I
am wrong. Perhaps it was a fluke of some sort that allowed Rubik (as contrasted
with Ishige, for example) to get most of the credit. But whatever the case, it does
illustrate my belief that people are extremely eager to attribute credit-even glory-to
just one




On Crossing the Rubicon                                                              357

person, and to vastly simplify a historical situation in order to be able to label it
and classify it in their minds.
        Who is willing to take the trouble to sift through all the murk surrounding
such monumental discoveries as relativity, Gödel’s incompleteness theorem,
digital computers, lasers, pulsars, the cosmic background radiation, or the structure
of DNA? Who wants to track down all the complexly tangled threads of ideas that
somehow led to one or two people getting all the glory? Almost without exception,
if you dig deep, you will find that the way the credit is conventionally apportioned
is unfair. Sometimes entirely the wrong person gets all the credit, sometimes
several unknown people deserve to share the credit, and sometimes the story is
even more complex and twisty than that. Somebody should write a book on bizarre
cases of credit attribution!
        But my point is simply that with the cube, as with anything that has made a
big hit, the world sees but the very tip of the iceberg, and someone in my position,
who receives a lot of mail on these matters, sees only a bit below the tip. There is a
lot more buried out there, and I am likely to get more letters from people who,
upon reading my current attempt to be fair (in other words, this Post Scriptum),
will feel especially slighted, given that I am trying to be fair and yet somehow
failed to mention their names! Ah, me, what can you do?

                                      *   *   *

        In the intervening time, I have not heard of any faster algorithm for solving
the Cube than Morwen Thistlethwaite's (described in Chapter 14). His algorithm,
originally known to solve the Cube in at most 52 turns, has now been slightly
improved on, thanks to computer searches. It is now known that 50 turns always
suffice, confirming a conjecture that Thistlethwaite himself had made several years
ago.
        Although this improvement on Thistlethwaite's algorithm does not
necessarily bring us appreciably closer to God's algorithm for the full 3 X 3 X 3
Cube, God's algorithm is now known for two important smaller puzzles: the 2 X 2
X 2 cube and Meffert's Pyraminx. Curiously, both of them require the same
number of turns at worst: eleven (disregarding the trivial turns of the tips of the
Pyraminx). The distribution of positions according to their distance from START
is quite interesting. Here it is for the Pyraminx, as supplied to me by John Francis
of Nutmeg, New Hampshire and Louis Robichaud of Sainte Foy, Quebec:

       1 configuration requires 0 moves
       8 configurations require 1 move
       48 configurations require 2 moves
       288 configurations require 3 moves
       1,728 configurations require 4 moves
       9,896 configurations require 5 moves




On Crossing the Rubicon                                                              358

       51,808 configurations require 6 moves
       220,111 configurations require 7 moves
       480,467 configurations require 8 moves
       166,276 configurations require 9 moves
       2,457 configurations require 10 moves
       32 configurations require 11 moves

Thus if START is at the "North Pole" of the space of all Pyraminx states, there are
32 different "South Poles", all maximally distant from it, and by far the bulk of the
population lives below the equator.
        By contrast, the 2 X 2 X 2 cube has 2,644 states at maximal distance
(eleven) from START. (In this metric, R2 counts as just one move, rather than
two.) Just as with the Pyraminx, the typical distance to START tends to be close to
the maximum distance, but that tendency is exaggerated even more in the 2 X 2 X
2. In particular, more than half the scrambled states require at least nine turns-and
yet, ten turns will suffice to reach over 99.9 percent of all states! Here is the
corresponding table:

       1 configuration requires 0 moves
       9 configurations require 1 move
       54 configurations require 2 moves
       321 configurations require 3 moves
       1,847 configurations require 4 moves
       9,992 configurations require 5 moves
       50,136 configurations require 6 moves
       227,536 configurations require 7 moves
       870,072 configurations require 8 moves
       1,887,748 configurations require 9 moves
       623,800 configurations require 10 moves
       2,644 configurations require 11 moves.

This information comes from the autumn-winter 1982 double issue of Singmaster's
Cubic Circular, and was apparently computed in several places around the world.

                                      * *    *

       In Chapter 14,1 described the game of inverting a handful of twists made
on a pristine Cube, and mentioned that Kate Fried could regularly invert seven and
once had undone ten. Peter Suber (the inventor of Nomic-see Chapter 4) calls this
challenge the "inductive game", and has mastered it to the same level as Kate Fried
did. He has written a short article describing this art, called "Introduction to the
Inductive Game of Rubik's Cube". In it he explains why he calls it that:

       The normal game is inductive only in the process a player undergoes in
discovering the algorithms sufficient for solution. That process has been said




On Crossing the Rubicon                                                             359

to model the scientific method, complete with the formulation and testing of
theories, negative results, and confirmation. The `inductive game' is inductive in
that way and more. The process of discovering the rules of mastery is similarly
inductive; but the product is also inductive. Instead of producing algorithms that
may be applied infallibly by an idiot, the inductive game produces `soft rules' or
probabilistic guides that must be applied in each case with judgment, mother wit,
and the weight of one's inductive experience ....
        The inductive game cannot become routine or boring, except perhaps to
gods. When one can solve three-twist randomizations nearly 100 percent of the
time, then one may move on to four-twist randomizations. Difficulty increases
exponentially. There is a foreseeable end to the series, of course. Players who
patiently gather up their nuanced, ineffable knowledge of random patterns may
reach 22-twist randomizations. Improvement does not merely approach the banal
satisfaction of more frequent success; it approaches hard knowledge of God's
algorithm.

       In the rest of his article, Suber details the results of his researches into this
domain and comes up with many hints and heuristics based on his notion of
information, defined as: "the adjacency of two or more tiles of the same color that
need not (and ought not) be separated on the shortest path home". His basic
guidelines (not to be interpreted overly rigidly) are:

       (1) Thou shalt not break up information.
       (2) Thou shalt endeavor to make more information.

        The catch is that many configurations give a false impression of containing
information. Suber calls this apparent information as distinguished from actual
information, and a large part of his article is devoted to hints for telling the two
apart. Readers interested in obtaining a copy of his article may write to Suber at the
Department of Philosophy, Earlham College, Richmond, Indiana 47374.

                                       *   *   *

        There is something tantalizing about the idea of precisely reversing a
scrambling. Suppose you could undo any scrambled state, and that one time the
resulting twist-sequence was found to be, say, UR-'D2LBLDR-'F2ULD-'BR-'U2L-
'DF. Would you be able to take this sequence apart and see any comprehensible
structure there? That is, would there be some recognizable pieces inside it that
explained how it undid that particular configuration?
        Another way of asking the same question is perhaps more compelling.
Some of my standard operators for getting things done on the Cube have the form
of commutators, conjugates, powers, or combinations of such things. For such
operators, I pretty much understand why they flip edges or do whatever they do.
However, there are a couple of operators in my repertoire that I've simply
memorized without having any understanding of




On Crossing the Rubicon                                                                360

why they accomplish what they do. For example, could it be that there's simply no
explanation why R -'D-'RD 'R -'D2R undoes three quarks on the bottom layer?
Could it be that there is, in other words, no conceptual breakdown to this operator?
Such a sequence would resemble a very, very large prime number, a structure that
admits of no breakdown into smaller "chunks".
        It seems almost certain that the shortest routes home from most scrambled
states on the Cube will admit of no breakdown; in short, that most of the solutions
given by God's algorithm are random, in the sense of having no internal rhyme or
reason to them-very much like a sequence of tosses of a coin or die. (This concept
of randomness is explained lucidly in the article "Randomness and Mathematical
Proof" by Gregory Chaitin.) If this is the case, it would mean that after a certain
point-most likely not far above ten twists-it will be a vain hope to try to undo a
Cube state via the route that got you there.
        Getting into a scrambled state and getting out of it are operations of
different computational complexity, just as getting yourself into a tight parking
space and getting yourself out of it are operations of different automobilistic
complexity. It is easier to find routes out than routes in, even though there are the
same number of each. (In this analogy, being well parked is the analogue of getting
to START, and being out in the street is the analogue of being scrambled.) Clearly,
there is something deeply asymmetric about such a situation, and the whole thing
smells of the second law of thermodynamics, stating that entropy will tend to
increase with time in a closed system.
        These informal intuitions can be made somewhat more precise. George
Marx, Eva Gajzago, and Peter Gnadig of the Department of Atomic Physics,
Eotvbs University, in Budapest, Hungary have studied the Cube statistically in a
paper called "The Universe of Rubik's Cube". To begin with, they define a face's
"color vector" as an ordered set of six numbers, telling how many facelets on that
face are red, orange, yellow, green, blue, and white, respectively. In START, the
red face's color vector is thus < 9,0,0,0,0,0 > . After some scrambling, you will get
color vectors more like this: <2,0,I,3,1,2>. Various numerical measures of any
face's "degree of scrambledness" can be derived from its color vector. The choice
made by these authors is the "length" of this vector-that is, the square root of the
sum of the squares of its "sides". For <9,0,0,0,0,0>, that comes out as 9, while for
the more typical < 2,0,1,3,1,2 >, it is about 4.36. The shortest possible color vector
consists of three l's and three 2's, and has length just under 4.
        Marx, Gajzago, and Gnadig studied the statistics of this quantity as the
Cube was twisted randomly, and discovered that faces whose color vector has
length 4.36 are the most common. Shorter or longer color vectors are quite
infrequent. If you start out at length 9 (a pristine Cube), then with random twisting
the length tends to decrease quickly to a bit under 5, and




On Crossing the Rubicon                                                              361

then to fluctuate around that value. This observation is their empirical formulation
of the second law of thermodynamics, establishing an "arrow of time".
        In accordance with standard usage in statistical mechanics, they define the
        entropy of a Cube face's state as the logarithm of the number of states that
have the same macroscopic description-in this case, the same color vector
(allowing rearrangements, so that < 2,0,1,3,1,2 > would be considered the same as
< 2,1,3,2,0,1 >). Then they show that standard formulas that apply to entropy in
real-world cases also apply to this "Cubical entropy". In particular, they remark:
"The distribution of the colored squares on a mixed-up cube can be described in a
similar way to how Maxwell and Boltzmann described the distribution of energy in
the molecular chaos of a gas." At the conclusion of their article, Marx, Gajzago,
and Gnadig wax lyrical: "I honor the cube as the smallest non-trivial model of the
great physical universe. " (italics theirs). (I suppose that when three authors jointly
describe themselves as "I", it is a case of "the editorial `I"'.)

                                   *      *       *

        During the Cube's peak popularity, a large number of speed tournaments
were held around the world and eventually a world champion emerged. He is Minh
Thai, formerly of Viet Nam, now resident in the United States. His winning time
on a scrambled Cube was 22.95 seconds. His average time seems to hover around
24 seconds, ranging as far upwards as 25 once in a while. Which leads me to ask:
Shouldn't he perhaps have been named Minh Time?
        There were also tournaments for the 4 X 4 X 4 cube, and there the best
times I heard of were in the three-minute range. Uwe Meffert sent me a 5 X 5 X 5
cube, which I must confess I never dared to scramble. I wonder how long the world
champion would take on that ! Meffert once described to me his dream of a "Magic
Triathlon", in which participants would have to unscramble a trio of scrambled
solids-as I recall, the objects involved were the Pyraminx, the Impossi*Ball, and
the Megaminx (Meffert's revised name for his Pyraminx Magic Dodecahedron-see
Figure 15-2d). My choices for the solids involved would have been different, but I
liked the basic idea. I do not know if such an event ever took place.
        I see no reason why harder events could not be created, involving such
esoteric skills as manipulating an N-dimensional cube represented in a computer,
such as H. R. Kamack and T. R. Keane's Magic Tesseract. These two gentlemen,
implementors of a 3 X 3 X 3 X 3 hypercube on a home computer, not only solved
the "basic mathematical problem" for this horrendous pseudo-object, but also
calculated the size of the group for the 3 X 3 X ... x3=3 ''' hypercube, or what they
call a "Rubik N-tope". For N=5, the size of this group is (approximately) 7.017 X
10560-a number not to be sneezed at!




On Crossing the Rubicon                                                               362

        According to Singmaster, mathematicians Joe Buehler, Brad Jackson, and
Dave Sibley studied the 3"' hypercube as well, and came up with a general
algorithm for it, as well as various conservation laws for it. The even more general
case of the M X M X M X ... X M =M'v hypercube remains unsolved, but it
particularizes (along another conceptual dimension) to the M X M X M cube.
Professor Jack Eidswick of the Department of Mathematics and Statistics at the
University of Nebraska sent me an article that presents an algorithm for solving
any member of this family of three-dimensional cubes. It is based on elaborate
versions of some of the necessary operators described in Chapter 14, built out of
conjugates and commutators and the like. I hear that Robert Brooks of the
Mathematics Department of the University of Maryland also has worked out such
an algorithm.

                                      *       *      *

        I finally must confront the matter of the cube fad's fading. David
Singmaster's Cubic Circular is going under after Volume 8. Many thousands of
Megaminxes were melted down for their plastic. Uwe Meffert's puzzle club seems
to have been a flop. The Skewb and many other wonderful objects I described
never hit the stands. A few that did were almost immediately gone forever. So ...
have we seen the last of the Magic Cube? Are those cubes you bought going to be
collector's items? Well, I am always loath to predict the future, but in this case I
will make an exception. I am bullish on the cube. It seemed to seize the
imagination wherever it went. Despite the line concluding my second cube column,
the cubic fad finally did spill over into the Soviet Union.
        In my opinion, the world simply overdosed on cube-mania for a while. We
humans are now collectively sick of the cube, but our turned-off state won't last too
long-no more than it lasts when you tell yourself "I'll never eat spaghetti again!"
after gorging on it. I predict that cubes will resurface slowly, here and there, and I
am even hopeful that some new varieties will appear now and then. This is Mother
Lode country. There may never again be quite the Gold Rush that we witnessed a
couple of years ago, but there's still plenty of gold in them thar hills!




On Crossing the Rubicon                                                              363

                          Mathematical Chaos
                         and Strange Attractors
                                   November, 1981

                                     I'm can't know how happy I am that we met,
                                     I'm strangely attracted to you.
                                                 -Cole Porter, "It's All Right with Me"


A few months ago, while walking through the corridors of the physics department
of the University of Chicago with a friend, I spotted a poster announcing an
international symposium titled "Strange Attractors". My eye could not help but be
strangely attracted by this odd term, and I asked my friend what it was all about. He
said it was a hot topic in theoretical physics these days. As he described it to me, it
sounded quite wonderful and mysterious.
        I gathered that the basic idea hinges on looking at what might be called
"mathematical feedback loops": expressions whose output can be fed back into them
as new input, the way a loudspeaker's sounds can cycle back into a microphone and
come out again. From the simplest of such loops, it seemed, both stable patterns and
chaotic patterns (if this is not a contradiction in terms!) could emerge. The difference
was merely in the value of a single parameter. Very small changes in the value of this
parameter could make all the difference in the world as to the orderliness of the
behavior of the loopy system. This image of order melting smoothly into chaos, of
pattern dissolving gradually into randomness, was exciting to me.




Mathematical Chaos and Strange Attractors                                           364

        Moreover, it seemed that some unexpected "universal" features of the
transition into chaos had recently been unearthed, features that depended solely on the
presence of feedback and that were virtually insensitive to other details of the system.
This generality was important, because any mathematical model featuring a gradual
approach to chaotic behavior might provide a key insight into the onset of turbulence
in all kinds of physical systems. Turbulence, in contrast to most phenomena
successfully understood in physics, is a nonlinear phenomenon: two solutions to the
equations of turbulence do not add up to a new solution. Nonlinear mathematical
phenomena are much less well understood than linear ones, which is why a good
mathematical description of turbulence has eluded physicists for a long time, and
would be a fundamental breakthrough.
        When I later began to read about these ideas, I found out that they had actually
grown out of many disciplines simultaneously. Pure mathematicians had begun
studying the iteration of nonlinear systems by using computers. Theoretical
meteorologists and population geneticists, as well as theoretical physicists studying
such diverse things as fluids, lasers, and planetary orbits, had independently come up
with similar nonlinear mathematical models featuring chaos-pregnant feedback loops
and had studied their properties, each group finding some quirks that the others had
not found. Moreover, not only theorists but also experimentalists from these widely
separated disciplines had simultaneously observed chaotic phenomena that share
certain basic patterns. I soon saw that the simplicity of the underlying ideas gives
them an elegance that, in my opinion, rivals that of some of the best of classical
mathematics. Indeed, there is an eighteenth- or nineteenth-century flavor to some of
this work that is refreshingly concrete in this era of staggering abstraction.
        Probably the main reason these ideas are only now being discovered is that the
style of exploration is entirely modern: it is a kind of experimental mathematics, in
which the digital computer plays the role of Magellan's ship, the astronomer's
telescope, and the physicist's accelerator. Just as ships, telescopes, and accelerators
must be ever larger, more powerful, and more expensive in order to probe ever more
hidden regions of nature, so one would need computers of ever greater size, speed,
and accuracy in order to explore the remoter regions of mathematical space. By the
same token, just as there was a golden era of exploration by ship and of discoveries
made with telescopes and accelerators, characterized by a peak in the ratio of new
secrets uncovered to money spent, so one would expect there to be a golden era in the
experimental mathematics of these models of chaos. Perhaps this era has already
occurred, or perhaps it is occurring right now. And perhaps after it, we will witness a
flurry of theoretical work to back up these experimental discoveries.
        In any case, it is a curious and delightful brand of mathematics that is being
done. This way of doing mathematics builds powerful visual imagery and intuitions
directly into one's understanding. The power of computers




Mathematical Chaos and Strange Attractors                                           365

allows one to bypass the traditional "theorem-proof-theorem-proof" brand of
mathematics, and to arrive quickly at empirical observations and discoveries that
reinforce each other, and that form a rich and coherent network of results. In the long
run, it may turn out to be easier to find proofs of these results (if proofs are still
desired), thanks to the careful and thorough exploration of the conceptual territory in
advance. It's an upstart's way of doing mathematics, and not all mathematicians
approve.
        One of the strongest proponents of this style of mathematizing has been
Stanislaw M. Ulam, who, when computers were still young, turned them loose on
problems of nonlinear iteration as well as on problems from many other branches of
mathematics. It is from Ulam's early studies with Paul Stein that many of the ideas to
be sketched here follow.

                                       *    *   *

         So much for romance. Let us work our way up to the concept of "strange
attractors" by beginning with the more basic concept of an attractor. This whole field
is founded on one concept: the iteration of a real-valued mathematical function-that is,
the behavior of the sequence of values x, f(x), f(f(x)), f(f(f(x))), ... , where f is some
interesting function. The initial value of x is called the seed. The idea is to feed f's
output back into f as new input over and over again, to see if some kind of pattern
emerges.
         An interesting and not too difficult problem concerning the iteration of a
function is this: Can you invent a function p with the property that for any real value
of x, p(x) is also real, and where p(p(x)) equals -x? The condition that p(x) be real is
what gives the problem a twist; otherwise the function p(x) = ix (where i is the square
root of - 1) would work. In fact, you can even think of the challenge as that of finding
a real-valued "square root of the minus sign". A related problem is to find a real-
valued function q, whose property is that q (q (x)) =1 /x for all x other than zero. Note
that no matter how you construct p and q, each will have the property that, given any
seed, repeated iteration creates a cycle of length four.
         Now, more generally, what kinds of functions, when repeatedly iterated, are
likely to exhibit interesting cyclic or near-cyclic behavior? A simple function such as
3x or x3, when iterated, does not do anything like that. The n th iteration of 3x, for
example, is 3 X 3 X 3 X ... X 3x, with n 3's-that is, 3"x-and the nth iteration of x3 is
just (((x3)3)3)"' s with n 3's again, which amounts to x3' . Nothing cycle-like here; the
values just keep going up and up and up. To reverse this trend, one needs a function
with some sort of switchback-a little zigzag or twist. A more technical way of putting
it is that one needs a nonmonotonic function: a function whose graph is folded-that is,
it starts moving one way-say upward-and then bends back the other way-say
downward.




Mathematical Chaos and Strange Attractors                                             366

      FIGURE 16-1. Two nonmomotonic, or “folded”, functions in the unit square.
      In (a) , a sharp peak, and in (b), a parabola. The maximum height of both is
      defined by the parameter λ.




Mathematical Chaos and Strange Attractors                                     367

              In Figure 16- la, we have a sawtooth with a sharp point at its top, and
      in Figure 16-1 b, a smoothly bending parabolic arc. Each of them rises from
      the origin, eventually reaches a peak height called X, and then comes back
      down for a landing on the far side of the interval. Of course there are
      uncountably many shapes that rise to height X and then come back down, but
      these two are among the simplest. And of the two, the parabola is perhaps the
      simpler, or at least the more mathematically appealing. Its equation is y
      =4Ax(1-x), with X not exceeding 1.
              We allow input (values of x) only between 0 and 1. As the graph
      shows, for any x in that interval, the output (y) always is between 0 and X.
      Therefore the output value can always be fed back into the function as input,
      which ensures that repeated iteration will always be possible. When you
      repeatedly iterate a "folded" function like this, the successive y- values you
      produce will sometimes go up and sometimes down-always hovering, of
      course, between 0 and X. The fold in the graph guarantees interesting effects
      when the function is iterated-as we shall see.
              It turns out that the spectacular differences in the degree of regularity
      of patterns I mentioned above are due to variations in the setting of what we
      might call the "X-knob". Depending on the value the knob is set at, the
      function yields an incredible variety of "orbits"-that is, sequences x,fix),
      f(f(x)), and so on. In particular, for X below ,a certain critical value
      X,=0.892486417967..., the orbits are all regular and patterned (although there
      are various degrees of patternedness; generally the lower it is, the more simply
      the orbit is patterned), but for X at or beyond this critical value, hold onto your
      hat! An essentially chaotic sequence of values will be traced out by the values
      x, f(x), f(f(x)),..., no matter what positive seed value of x you choose. In the
      case of the parabola, the critical role played by varying the A-knob seems to
      have been first realized by P. J. Myrberg in the early 1960's, but his work was
      published in a little-known journal and did not attract much attention. Some
      ten years later, Nicholas C. Metropolis, Paul Stein, and Myron Stein
      rediscovered the importance of the knob not only for the parabola but also for
      many other functions. Indeed, they discovered that as far as certain topological
      properties were concerned, the function did snot matter-only the value of A
      did. This property has come to be called "structural universality".

                                          *    *   *

              In order to see how such a nonintuitive dependence on the setting of
      the it-knob comes about, one must develop a visual sense for the process of
      iterating f(x). This is readily done. Suppose we set it to 0.7. The graph of f(x)
      appears in Figure 16-2. In addition, the line y =x appears as, a 45-degree
      broken line. (This graph and most of the others in this article were produced
      on a small computer by Mitchell J. Feigenbaum of the Los Alamos National
      Laboratory.)




Mathematical Chaos and Strange Attractors                                            368

      FIGURE 16-2. The parabola defined by -X-knob" setting of 0.7. An initial x-
      value of about 0.04 is used as a "seed "for iteration, and the pathway taken is
      shown. Eventually it settles down at a fixed point, denoted by x*.

              Consider the two x-values where the 45-degree line and the curve
      intersect. They are at x=0 and x=9/14=0.643. Let us designate the nonzero
      value as x'. By construction, then, f(x') equals x', and repeated iteration off at
      this x-value will get you into an infinite loop. The same happens if you start
      iterating at x=0: you get stuck in an endless loop. However, there is a
      significant difference between these two fixed points of f It is best indicated by
      taking some other initial value of x, say one close to 0.04, as is shown in the
      same figure. Call this starting x-value x 0 . There is an elegant graphical way to
      generate the orbit of any seed x 0 . A vertical line at x-value xo will hit the
      curve at height y0 =f(x 0 ). To iterate f, we must draw a new vertical line located
      at the new x-value equal to this y-value. This is where the 45-degree line y =x
      comes in handy. Staying at height y0 , we simply move over horizontally until
      we hit that 45-degree line. Then, since along this line y equals x, both x and y
      equal y 0 . Let us call this new x-value x,. We now draw a second vertical line.
      This one will hit the curve at height y1 = f(x 1 )=f(y 0 )=f(f(x 0 )). Now we just
      repeat.




Mathematical Chaos and Strange Attractors                                            369

In brief, iteration is realized graphically by a simple recipe:

       (1) Move vertically until you hit the curve; then
       (2) Move horizontally until you hit the diagonal line.
       Repeat steps (1) and (2) over and over again.

         The results of this procedure with seed x0=0.04 are also shown in Figure 16-2.
We are led in a merry chase 'round and 'round the point whose x-coordinate and y-
coordinate are x'. Gradually we close down on that point. Thus x' is a special kind of
fixed point, because it attracts iterated values of f(x). It is the simplest example of an
attractor: every possible seed (except 0) is drawn, through iteration off, to this stable
x-value. This x' is therefore called an attractive or stable fixed point. By contrast, 0 is
a repellent or unstable fixed point, since the orbit of any initial x-value, even one
infinitesimally removed from 0, will proceed to move away from 0 and toward x'.
Note that sometimes the iterates off will overshoot x', sometimes they will fall short-
but they inexorably draw closer to x', zeroing in on it like swallows returning to
Capistrano. One might also think of such familiar and charming metaphors of prey-
seekers as heat-seeking missiles, mosquitos, bloodhounds, Nazi-hunters, sharks, and
lastly, the children's rhyme, "Around the world, and around the world, goes a big
bear; he bores a hole, and he bores a hole, right ... in . . . there!"
         What accounts for this radical qualitative difference between the two fixed
points (0 and x') off? A careful look at Figure 16-2 will show that it is the fact that at
0, the curve is sloped too steeply. In particular, the slope there is greater than 45
degrees. It is the local slope of the curve that controls how far you move horizontally
each time you iterate f. Whenever the curve is steeper than 45 degrees (either rising or
falling) it tends to pull you farther and farther away from your starting point as you
repeatedly iterate by rules (1) and (2). Hence the criterion for the stability of a fixed
point is: The slope at the fixed point should be less than 45 degrees. Now, this is the
case for x' when A equals 0.7. In fact, the slope there is about 41 degrees, whereas at
0 it is much greater than 45 degrees.
         What happens if we increase A? The position of x' (x' being by definition the
x-coordinate of the point where the curve f and the line y =x intersect) will change,
and the slope off at x' will increase as well. What happens when the slope hits 45
degrees or exceeds it? This occurs when A is 3/4. We will call this special value of the
A-knob A,. Let us look at the graph for a slightly greater A-knob setting, namely
A=0.785. (See Figure 16-3.)
         What if we begin with some random seed instead, again say x = 0.04? The
resulting orbit is shown in Figure 16-3a. As you can see, a very pretty thing happens.
At first the values move up toward the vicinity of x' (now an unstable fixed point of
f), but then they spiral gradually outward and settle




Mathematical Chaos and Strange Attractors                                              370

FIGURE 16-3. Here, the λ-knob has box twisted a little higher: now it is set at 0.785.
On top, f itself is shown, and below, f’s two-humped iterate g is shown. Now a typical
seed will eventually settle into a two-step square dance, asymptotically jumping back
and forth between values x 1 * and x 2 *




Mathematical Chaos and Strange Attractors                                         371

down smoothly to a kind of "square dance" converging on two special values x, and
x2. This elegant oscillation is called a 2-cycle, and the pair of x-values that constitute
it (x; and x2) is again an attractor-in particular, an attractor of period two. This term
means that our 2-cycle is stable: it attracts x-values from far and wide as f is iterated.
The orbit for any positive seed value (except x' itself) will eventually fall into the
same dance. That is, it will asymptotically approach the perfect 2-cycle composed of
the points x,` and X2\% although it will never quite reach it exactly. From a physicist's
point of view, however, the accuracy of the approach soon becomes so great that one
can just as well say that the orbits have been "trapped" by the attractor.
        An enlightening way to understand this is to look at a graph of a new function
made from. the old one. Consider the graph of g (x) =f(f(x)), shown in Figure 16-3b.
This two-humped camel is called the iterate off. First of all, observe that any fixed
point off is also a fixed point of g, so that 0 and x' will be fixed points of g. But
secondly observe that since f(x,) equals x2, and conversely f(x2) equals x,, g will have
two new fixed points: g(x,) =x; and g(x2)=x2. Graphically, x; and x2` are easily
found: they are intersection points of the 45-degree line with the two-humped graph
of g (x). There are four such points (0 and x' being the other two). As we have seen,
the criterion for the stability of any fixed point under iteration is that the slope at that
point should be less than 45 degrees. Here we are concerned with fixed points of g,
and hence with g's slope (as distinguished from f's slope). Indeed, in the same figure,
you can clearly see that at 0 and at x', g is sloped more steeply than 45 degrees,
whereas at both x; and x2, g's slope is less than 45 degrees. In fact, quite remarkably,
not only are both slope values less than 45. degrees, but also, as it turns out through a
simple bit of calculus, they are equal (or "slaved" to each other, as it is sometimes
put).

                                         *   *   *

          We have now seen an attractor of period one get converted into an attractor of
period two at a special value of λ (namely, λ =3/4). Precisely at that value, the single
fixed point x' splits into two oscillating values, x, and x2. Of course they coincide at
"birth", but as λ increases, they separate and draw farther and farther apart. This
increase of λ will also cause g's slope at these two stable fixed points (of g) to get
steeper and steeper until finally, at some λ -value, g, like its progenitor f, will reach its
own breaking point (i.e., the identical slopes at both x; and x2 will exceed 45
degrees), and each of these two attracting points will break up, spawning its own local
2-cycle. (Actually, the cycles are 2-cycles only as far as g is concerned; for f, the new
points are elements of an attractor of period four. You must be careful to keep f and g
straight in your mind!) These two splittings will happen at exactly the same "moment"
(i.e., at the same λ -knob setting), since the value of the slope of g at x; is slaved to the
value of the slope at x2. This λ -knob setting will be called λ 2, and it has the value of
0.86237 ...




Mathematical Chaos and Strange Attractors                                               372

FIGURE 16-4. A picture of f's iterate's iterate h at a still higher value of λ, namely
0.87.

         Here, as with a joke, you may anticipate the punch line by the time you have
heard the theme and one variation. Hence by now you have probably surmised that at
some new value λ, all four points in f's attractor will simultaneously fission, yielding a
periodic attractor consisting of eight points; and thereafter this pattern will go on and
on, doubling and redoubling as various special λ-knob settings are reached and
passed. If this is your guess, you are quite right, and the underlying reason is the same
each time: the (identical) slopes at all the stable fixed points of some graph reach the
critical angle of 45 degrees. In the case of the first fission (at λ,) it was the slope off
itself at the single point x'. The next fission was due to the slopes at g's two stable
fixed points x; and x2 simultaneously reaching 45 degrees. Analogously, A3 is that
value of A at which the slope of h (x) = g(g(x))=f(f(f(f(x)))) hits 45 degrees
simultaneously at the four stable fixed points of h. And so it goes. Figure 16-4 shows
the bumpy appearance of h (x) at a A-value of approximately 0.87.
         In Figure 16-5, the locations on the x-axis of the stable fixed points of f are
shown for A, through A6 (by which time there are 32 of them, some clustered so
closely that they cannot be distinguished). The points are pictured just at the moment
of their becoming unstable, each one like a cell




Mathematical Chaos and Strange Attractors                                              373

FIGURE 16-5. Showing how stable attractors become unstable and undergo `fusion "
at a series of increasing A-values, denoted A„ for n=1, 2, 3,... Note how the boxed
subpattern on the lowest line resembles the entire pattern two lines above. This
resemblance becomes more and more accurate the larger n gets.

FIGURE 16-6. A graph showing the evolution of attractors as A increases from 0 to
1. Bifurcations begin at A=0.75 and escalate towards chaos. The "chaotic region',
beginning at A= 0.892 shows unexpectedly beautiful fine structure. [From "Roads to
Chaos " by Leo P. Kadanof in Physics Today, December 1983 p.51; see also J. P.
Crutchfield, D. Farmer, and B. A. Huberman, Physics Reports, Vol. 92, pp. 45-82,
December, 1982. )




Mathematical Chaos and Strange Attractors                                      374

on the verge of division. Notice the neat pattern in the distribution of the attracting
points. Looking at these graphs of the spacings of the elements of the successive
period-doubled attractors off, you can see that each line can be made from the one
above it through a recursive geometric scheme whereby each point is replaced by two
"twin" points below it. Each local clustering pattern of points echoes the global
clustering pattern, simply reduced in scale (and also, in alternating local clusters, left
and right are reversed). For example, in the bottom line a local group of eight points
has been outlined in color. Notice how the group of points is like a miniature version
of the global pattern two lines above it.
         The discovery of this recursive regularity, first made on a little calculator by
Feigenbaum, is one of the major recent advances in the field. It states in particular that
to make line n + 1 from line n, you simply let each point on line n give birth to
"twins". The new generation of points should be packed in about 2.5 times more
densely than the old generation was. More exactly .stated, the distance between new
twins should be a times smaller than the distance between their parent and its twin,
where a is a constant, approximately equal to 2.5029078750958928485 ... This rule
holds with greater and greater accuracy the larger n becomes.
         What about the values of the A's? Are they headed asymptotically toward 1 ?
Surprisingly enough, no. These A-values are quickly converging on a particular
critical value A, of size roughly 0.892486418 ... And their convergence is remarkably
smooth, in the sense that the distance between successive A's is shrinking
geometrically. More precisely, the ratio (A,-A,-,)/(A,+,-A.) approaches a constant
value called 8 by Feigenbaum, its discoverer, but more often referred to simply as
"Feigenbaum's        number"      by     others.    Its     value    is    approximately
4.66920160910299097...
         In short, as λ approaches λ c , at special A-values predicted by Feigenbaum's
constant 8, f's attractor doubles in population, and its increasingly many elements are
geometrically arranged on the x-axis according to a simple recursive plan, the main
determining parameter of which is Feigenbaum's other constant, a.
         Then for λ beyond λ c -called the chaotic regime-the results of iterating f can,
for some seed values, yield orbits that converge to no finite attractor. These are
aperiodic orbits. For most seed values, the orbit will remain periodic, but the
periodicity will be very hard to detect. First of all, the period will be extremely high.
Secondly, the orbit will be much more chaotic than before. A typical periodic orbit,
instead of quickly converging to a geometrically simple attractor, will meander all
over the interval [0,1 ], and its behavior will appear indistinguishable from total
chaos. Such behavior is termed ergodic. Furthermore, neighboring seeds may, within
a very small number of iterations, give rise to utterly different orbits. In short, a
statistical view of the phenomena becomes considerably more reasonable beyond λ c .
         Figure 16-6 beautifully portrays the period-doubling route to chaos, as




Mathematical Chaos and Strange Attractors                                             375

well as what happens after you've gotten there. The bifurcations are clear to the eye,
and since the horizontal distance from each set of them to the next shrinks
geometrically, the onset of chaos at λ is plainly visible. But the regularity of the
structure to the right of λ c ,-that is, in the chaotic regime -is quite unexpected. It is
certain that there are many deep mathematical secrets locked up in this elegant graph.

                                        *   *   *

        Now, what do such novel concepts as the iteration of folded functions, period
doubling, chaotic regime, and so on have to do with the study of turbulence in
hydrodynamic flow, the erratic population fluctuations in predator-prey relations, and
the instability of laser modes? The basic idea is embedded in the contrast between
laminar flow and turbulent flow. In a peacefully flowing fluid, the flow is laminar-a
soft and gentle word that means that all the molecules in the fluid are moving like cars
on a multilane freeway. The key features are: (1) that each car follows the same path
as its predecessor, and (2) that two nearby cars, whether they are in the same lane or
in different ones, will, as time passes, slowly separate from each other essentially in
proportion to the difference in their velocities-which is tb say, linearly. These features
also apply to molecules of fluid in laminar flow; there, the lanes are called streamlines
or laminas.
        By contrast, when a fluid is churned up by some external force, this smooth
behavior turns into turbulent behavior, as is seen in breakers at the beach and cream
being stirred into coffee. Even the word "turbulent" .sounds much harsher and more
angular than the soft word "laminar". Here, the image of a multilane freeway no
longer holds; the streamlines separate from each other and tangle in the most
convoluted of ways, as shown in Figure 16-7. In such systems there are eddies and
vortices and all sorts of unnamable whorls on many size-scales at once, and
consequently, two points that were initially very close may soon wind up in totally
different regions of the fluid. Such quickly diverging paths are the hallmark of
turbulence. The distance between points can increase exponentially with time, instead
of just linearly, and the coefficient of time in the exponent is called the Lyapunov
number. When one speaks of chaos in turbulent flow, it is this rapid, nearly
unpredictable separation of neighbors that is meant. Such behavior is strikingly
reminiscent of the rapid separation, in the chaotic regime of X, of two orbits whose
seeds might originally have been very close together.




Mathematical Chaos and Strange Attractors                                             376

FIGURE 16-7. Showing the approach to turbulence. In the upper two pictures, a rod
was drawn through a viscous liquid once, setting up trains of vortices behind it. In the
lower two, the rod was drawn more than once, and the forms are therefore more
complicated and recursive seeming. It is provocative to compare this figure with
Figure 13-4. [From Sensitive Chaos, by Theodor Schwenk.




       )


Mathematical Chaos and Strange Attractors                                           377

This suggests that the "scenario" (as it is called) by which pretty, periodic orbits
gradually give way to the messy, chaotic orbits of our parabolic function might
conceivably be mathematically identical to the scenario underlying the transition to
turbulence in a fluid or other system. Exactly how this connection is established,
though, requires some more detailed setting of context. In particular, we must briefly
consider how the spatio-temporal flow of a fluid or some other entity, such as
population density or money, is mathematically modeled.
        In such real-world problems, the most successful equations yet found to model
the phenomena are differential equations. A differential equation connects the
continuous rate of variation of some quantity to that quantity's current size and the
current sizes of other quantities. Moreover, the time variable is itself continuous, not
jerking from one discrete instant to the next as some strange clocks and watches
occasionally do, but indivisibly flowing, like a liquid. One way to visualize the
patterns defined by differential equations is to imagine a multidimensional space-it
could have thousands of dimensions, or merely a few-in which a point is continuously
tracing out a curve. At any one moment, the single point contains all the information
about the state of the physical system. Its projections along the various axes give the
values of all the relevant quantities that pin down a unique state. Clearly the space-
called phase space-would need to have an enormous number of dimensions for a mere
point to store the entire shape of a wave breaking on a beach. On the other hand, in a
simple predator-prey relation, only two dimensions suffice: one coordinate, say x,
giving the predator population and the other, say y, giving the prey population. Two
dimensions are more easily visualized, and so we will stick with that case for the time
being. The ideas generalize, however, to higher-dimensional cases.
        As time progresses, x and y determine each other in an intertwined manner.
For example, a large population of predators will tend to reduce the population of
prey, whereas a small population of prey will tend to reduce the population of
predators. In such a system, x and y constitute a single point (x,y) that swirls around
smoothly in a continuous orbit on the plane. (Here the sense of "orbit" is different
from the preceding one-that of the discrete, or jumping, orbits we saw when our
parabolic function was iterated.) One such possible orbit appears in Figure 16-8; it is
generated by a differential equation called "Duffing's equation". It looks like the path
of a buzzing fly in your bedroom-or rather, it looks like the shadow of the fly's path
on a wall. As a matter of fact, this self-intersecting two-dimensional curve is the
shadow of a non-self-intersecting three-dimensional curve. The motion of a point in
phase space must always be non-self-intersecting. This arises from the fact that a
point in phase space representing the state of a system encodes all the information
about the system, including its future history, so that there cannot be two different
pathways leading out of one and the same point.
In particular, in Duffing's equation there is a third variable, z, that I have




Mathematical Chaos and Strange Attractors                                           378

FIGURE 16-8. If values of x and y mutually determine each other according to
Duffing's equation as time passes, then the point (x,y) will trace out a curve (a). If a
strobe light blinks periodically and shows (x,y) at selected instants, then a group of
isolated points will start appearing as in (b), and gradually filling out a region of
their own-a Poincare map.




Mathematical Chaos and Strange Attractors                                           379

not mentioned so far. If you think of x and y as representing predator and prey
populations, then you can think of z. as representing a periodically varying external
influence, such as the sun's azimuth or the amount of snow on the ground. Now, if
you will allow me to mix my buzzing-fly image with the predator-prey example,
imagine a bedroom with a fly buzzing periodically back and forth between two walls.
Let us say it takes the fly a year to cross the room and come back. (Perhaps it is a
rather large bedroom, or maybe just a slow fly.) In any case, as the fly flies, its
shadow on one of the two walls traces out the curve shown in Figure 16-8a. If the fly
ever chances to come back to a point in the room which it has passed through before,
it is doomed to loop forever, following the path it took the preceding time over and
over again. This gives you a picture of the continuous orbit of a point in phase space
representing the state of dynamic system controlled by differential equations.

                                         *   *   *

        Now suppose we wanted to establish some connection of these systems to
discrete orbits. How might we do so? Well, the values of x, y, and z need not be
watched at all moments; they can be sampled periodically, at some natural frequency.
In the case of animal populations, a year is the obvious natural period. The sun's
azimuth is exactly periodic, and the weather at least tries to repeat itself a year later.
Thus a natural sequence of discrete points (x,, y,, z, ), (x2, y2, z2), . . . can be singled
out-one per year. It is as if a strobe light blinked regularly and froze the fly on special
annual occasions-perhaps at midnight every Halloween. Or you can think of a firefly
that flashes on for just a split second once every year. At all other times its
peregrinations around the room are unseen. Figure 16-8b shows a sequence of discrete
points along the fly-path's shadow, marked by numbers telling when they occurred.
Gradually, as many "years" elapse, enough of these discrete points will accumulate
that they will start to form a, recognizable shape of their own. This pattern of points is
a discrete "orbit", and so it is closely related to the discrete orbits defined by the
iteration of our parabola f(x). In that parabolic case, we had a simple one-dimensional
recurrence relation (or an iteration):

       xn+l=f(xn)

Here we have a two-dimensional recurrence:

       xn+i =fl(xn,yn)

       Yn+I =f2(xn,yn)

This is a system of coupled recurrence relations, in which output values of the nth
generation (xn,yn)are fed right back into f, and f2 as new inputs, to




Mathematical Chaos and Strange Attractors                                              380

produce the n + 1st generation. On and on it goes, generation after generation. In
higher-dimensional cases, of course, there are more such '~ equations. Nevertheless,
the skeleton of all these systems remains the same: a multidimensional point (xn,yn,
z) jumps from one discrete location in phase space to another, as a discrete variable, n,
representing time jumping ahead in discrete units, is incremented.
        Notice that we have finessed our way around the continuous time variable that
is involved in differential equations. We have done it by focusing on the' way the
point is connected to its predecessor one "year" earlier (or whatever natural period is
involved). But is there always a "natural period" at which to look at a system of
mutually intertwined differential equations? Not always. In some situations, however,
there is, and this happens to be the case in all situations where turbulent behavior
occurs.
        Why is this so? All systems that exhibit turbulent behavior are dissipative,
which means that they dissipate, or degrade, energy from more usable forms such as
electricity into the less usable form of heat. In the case of hydrodynamic flow, this
dissipation is caused by friction, and in the other systems we have been considering,
by abstract analogues of friction. A familiar consequence of friction is that objects in
motion will grind to a halt unless energy is pumped in. Now if we "drive" a
dissipative system with a periodic driving force (you can imagine, for example,
stirring a cup of coffee with a spoon in a periodic, circular way), then, of course, the
system will not grind to a halt; it will head for some kind of steady state. Such a
steady state is a stable orbit-or in our terms, an attractor in phase space. And since we
have driven the system with a periodic spoon, we have defined a natural frequency at
which to flash our strobe light and freeze the system's statenamely, each time the
spoon comes swinging around and passes some fixed mark on the cup, such as its
handle. This will constitute our "year". In this way, continuous time can be replaced
by a series of discrete instants, as long as we are dealing with a dissipative system
driven by a periodic force. And so continuous orbits can be replaced by discrete
orbits, which brings iteration back into the picture.
        If the driving force itself has no natural period (it may be simply a constant
force), there is still a way to define a natural period, as long as some variable in the
system swings back and forth between extremes. Just flash your strobe whenever that
variable hits its extreme value, and the fly will still be caught at discrete instants. This
type of discrete representation of the fly's motion in a multidimensional space is
called a Poincare map.
        This stirring argument is only hand-waving, of course, and needs much more
rigor to be convincing to a mathematician. It nonetheless gives the flavor of how the
study of a set of coupled differential equations can be replaced by the study of a set of
coupled discrete recurrence relations. This is the vital step that brings us back to the
recent discoveries about the parabola.

                                         *   *   *




Mathematical Chaos and Strange Attractors                                              381

In 1975, Feigenbalum discovered that his numbers a and S do not depend on the
details of the shape of the curve defined by f(x). Almost any smooth convex shape
that peaks in the same spot will do as well. Inspired by the structural universality
discovered by Metropolis, Stein, and Stein, Feigenbaum tried working with a sine
curve instead of a parabola. He was flabbergasted by the reappearance of the same
numerical values, to many decimal places, of the numbers a and S, which had
characterized the period-doubling and the onset of chaos for the parabola. For the sine
curve just as for the parabola, there is a height-parameter A and a set of special A-
values that converge to a critical point Ac. Moreover, the onset of chaos at Ac is
governed by the same numbers a and S. Feigenbaum began to suspect that there was
something universal going on here. In other words, he suspected that what is more
important than f itself is the mere fact that f is being iterated over and over. In fact, he
suspected that f itself might play no role in the onset of chaos.
        It is not quite that simple, in reality. Feigenbaum soon discovered that what
does matter about f is just the nature of the peak at its very center. The long-term
behavior of orbits depends only on an infinitesimal segment at the crest of the graph,
and ultimately, it depends only on the behavior at the very point where the maximum
occurs! The rest of the shape, even the region close to the peak, is irrelevant. A
parabola has what is called a quadratic maximum, as do a sine wave, a circle, and an
ellipse. In fact, the behavior of a randomly-produced smooth function at a typical
maximum would be expected to be of the quadratic type, in the absence of any special
coincidences. So the parabolic case, rather than being a quirky exception, begins to
seem like the rule. This empirical discovery by Feigenbaum, involving two
fundamental scaling factors a and S that characterize the onset of chaos through
period-doubling attractors, represents a new kind of universality, known as metrical
universality, to distinguish it from the earlier-known structural universality. This
empirically demonstrated metrical universality was later proved to be correct (in the
more orthodox sense of proof) in the one-dimensional case by Oscar Lanford.
        A truly exciting development occurred when Feigenbaum's constants
unexpectedly turned up in some messy models of actual physical systems that exhibit
turbulence, not just in pretty and idealized mathematical systems. Valter Franceschini
of the University of Modena in Italy adapted the Navier-Stokes equation, which
governs all hydrodynamic flow, for computer simulation. To do so, he turned it into a
set of five coupled differential equations whose Poincare maps he could then study
numerically on his computer. He first found that the system exhibited attractors with
repeated period-doubling as its governing parameters approached the values where
turbulence was expected to set in. Unaware of Feigenbaum's work, he showed his
results to Jean-Pierre Eckmann of the University of Geneva, who immediately urged
him to go back and determine the rate of convergence of the A-values at which
period-doubling occurred. To their




Mathematical Chaos and Strange Attractors                                              382

amazement, Feigenbaum's a- and 8-values-accurate to about four decimal places-
appeared seemingly out of nowhere! For the first time, an accurate mathematical
model of true physical turbulence revealed that its structure was intimately related to
the humble chaos lurking in the humble parabola y=4Ax(1-x). Subsequently,
Eckmann, Pierre Collet, and H. Koch showed that in the behavior of a
multidimensional driven dissipative system, all 'dimensions but one tend to drop out
after a sufficiently long period of time, and so one should expect the characteristic of
one-dimensional behaviour namely Feigenbaum's metrical universality-to reappear.
         Since then, experimentalists have been keeping their eyes peeled for period-
doubling behavior in actual physical systems (not just in computer models). Such
behavior has been observed in certain types of convective flow, but so far the
measurements are too imprecise to lend very strong support to the idea that the
parabola contains the clues revealing the nature of genuine physical turbulence. Still,
it is tantalizing to think that somehow, all that really matters is that a dissipative set of
coupled recurrence relations is being iterated-but that the detailed properties of those
recurrences can be entirely ignored if one is concentrating on understanding the route
to turbulence.
         Feigenbaum puts it this way. One often sees a pattern of clouds in the sky -a
celestial trellis composed of a myriad of small white puffs stretching from horizon to
horizon-that clearly did not happen "by accident". Some systematic hydrodynamic
law has got to be operating. Yet, says Feigenbaum, it must be a law operating at a
higher level, or on a larger scale, than the Navier-Stokes equation, which is based on
infinitesimal volumes of fluid and not on large "chunks". It seems that in order to
understand such beautiful sky patterns, one must somehow bypass the details of the
Navier-Stokes equation, and come up with some coarser-grained but more relevant
way of analyzing hydrodynamic flow. The discovery that iteration gives rise to
universality-that is, independence of the details of the function (or functions) being
iterated-offers hope that such a view of hydrodynamics may be well on its way to
emerging.

                                     *       *       *

Well, we have covered attractors and turbulence; what about strange attractors? We
have now built up the necessary concepts to understand this idea. When a periodically
driven two-dimensional (or higher-dimensional) dissipative system is modeled by a
set of coupled iterations, the successive points lit up by the flashes of the periodic
strobe light trace out a shape that plays the role, for this system, that a simple orbit did
for our parabola. But when one is operating in a space of more than one dimension,
the possibilities are richer. Certainly it is possible to have a stable fixed point (an
attractor of period one). This would just mean that at every flash of the strobe, the
point representing the system's state is exactly where it was last




Mathematical Chaos and Strange Attractors                                               383

time. It is also possible to have a periodic attractor: one where after some finite
number of flashes, the point has returned to a preceding position. This would be
analogous to the 2-cycles, 4-cycles, and so on that we saw occurring for the parabola.
        But there is another option: that the point never returns to its original position
in phase space, and that successive flashes reveal it to be jumping around quite
erratically inside a restricted region of phase space. Over a period of time, this region
may take shape before an observer's eyes as the strobe flashes periodically. In the
majority of such cases so far studied, a most unexpected phenomenon has been
observed to take place: the erratically jumping point gradually creates a delicate
filigree that recalls the "faint fantastic tracery made by frost on glass". (I owe this
poetic image to the American critic James Huneker, who used it to describe the
magical effect of one of Chopin's piano etudes: Op. 25, No. 2-see Chapter 9.) The
delicacy is of a rather specific kind, closely related to th, "fractal" curves described by
Benoit Mandelbrot in his book The Fractal Geometry of Nature. In particular, any
section of such an attractor, when blown up, reveals itself to be just as exquisitely
detailed as was the larger picture from which it was taken. In other words, there is an
infinite regress of detail, a never-ending nesting of pattern within pattern. One of the
earliest of such structures to be found, called the attractor of Henon, is shown in
Figure 16-9. It is generated by the sequence of points (xn,y„) defined by the following
recurrence relations:

       xn+1 =Yn-ax„ 2-1

       Yn+l =bx„

Here, a is equal to 7/5 and b to 3/10; the seed values are x0=0 and y0=0. The small
square in Figure 16-9a is blown up in Figure 16-9b to reveal more detail, and then
another square in Figure 16-9b is blown up in Figure 16-9c to reveal yet finer detail.
Note that what we appear to have is a sort of three-lane highway each of whose lanes
breaks up, when magnified, into more parallel lanes, the outermost of which is a new
three-lane highwayand on and on it goes. Any perpendicular cross-section of this
highway would be what is called a "Cantor set", formed by a simple and famous
recursive process.
        Begin with a closed interval, say [0,1 ]. ("Closed" means that the interval
includes its endpoints.) Now eliminate some open central subinterval. (Since an open
subinterval does not include its endpoints, those two points will remain in the Cantor
set being constructed before your eyes.) Usually the deleted subinterval is chosen to
be the middle third (1/3, 2/3), but this is not necessary. Two closed subintervals
remain. Subject them to the same kind of process-namely, eliminate an open central
subinterval inside each of them. Repeat the process ad infinitum. What you will be
left with at the end of your infinite toil will be a delicate structure consisting of
isolated points stretched out along the original segment [0,1] like beads of dew on




Mathematical Chaos and Strange Attractors                                              384

FIGURE 16-9. The attractor of Henon: a strange attractor. In (a), the full curve is
shown. In (b), the boxed region of (a) is blown up to reveal hidden details. In (c), the
boxed region of (b) is further blown up to reveal yet more deeply hidden details. And
on and on it could go, ad infinitum.




Mathematical Chaos and Strange Attractors                                           385

a wire. Their number, however, will be uncountably infinite, and their density will
depend on the details of your recursive elimination process. Such is the nature of a
Cantor set, and if an attractor's cross-sections have this weird kind of distribution, the
attractor is said to be strange, and for good reason.
          Another beautiful strange attractor is generated by the "stroboscopic" points 0,
1, 2,... in Figure 16-8b. Since this pattern comes out of Duffing's equation, it is called
"Duffing's attractor", and it is shown in a slightly expanded scale in Figure 16-10.
Notice its remarkable similarity to the attractor of Henon. Could this be universality
showing its face again?
          It is interesting that for the parabola, at the critical value X,, f's attractor
suddenly consists of infinitely many points, since it is the culmination of an infinite
sequence of bifurcations. You can visualize this set either as the limiting case of the
horizontal point-sets in Figure 16-5, or as the vertical point-set belonging to x=X in
Figure 16-6. The precise scatter-pattern of this uncountable point-set is determined by
Feigenbaum's recursive rule involving his constant a. Given its recursive genesis, it
seems probable that this particular attractor is a Cantor set. Hence the fertile parabola
has provided us with an example of a one-dimensional strange attractor!
In the chaotic regime of the more general k-dimensional case, long-term prediction of
the path that a point will take is quite impossible. Two nearly

FIGURE 16-10. The strange attractor that emerges from a Poincare map of Duffing's
equation.




Mathematical Chaos and Strange Attractors                                             386

touching points on a strange attractor will, after a few blinks of the strobe light, have
wound up at totally different places. This is called sensitive dependence on initial
conditions and is another defining criterion of a strange attractor.

                                       *   *   *

         At present, no one knows just why, how, or when strange attractors will crop
up in the chaotic regimes of iterative schemes representing dissipative physical
systems, but they do seem to play a central role in the mystery of turbulence. David
Ruelle, one of the prime movers of this whole approach to turbulence, wrote: "These
systems of curves, these clouds of points, sometimes evoke galaxies or fireworks,
other times quite weird and disturbing blossomings. There is a whole world of forms
still to be explored, and harmonies still to be discovered."
         Robert M. May, a theoretical biologist, concluded his now quite famous
review article covering the field in 1976 with a plea that I find most apt and would
like to repeat:

          I would urge that people be introduced to the equation y =4 kx (1-x)
   early in their mathematical education. This equation can be studied
   phenomenologically by iterating it on a calculator, or even by hand. Its study
   does not involve as much conceptual sophistication as does elementary calculus.
   Such study would greatly enrich the student's intuition about nonlinear systems.
          Not only in research but also in the everyday world of politics and
   economics, we would all be better off if more people realized that simple
   nonlinear systems do not necessarily possess simple dynamical properties.

Post Scriptum.
Stanislaw Ulam, a uniquely inventive mathematician and a warm and delightful
human being, died as I was working on this series of postscripts. I had the good
fortune to get to know Stan Ulam and his French-born wife Francoise in the summer
of 1980, when I visited Santa Fe and stayed with them for a few days. I had always
admired and felt kinship with Ulam's strange style in mathematics, totally driven by a
passion for the the quirky and the unpredictable, bored by the pure and regular. Ulam
loved more than anything to find total chaos in the midst of pristine order. Of course,
the thrill was in knowing that there was some kind of law to this chaos, so that in
reality-that is, in God's eye-there was simply a deeper kind of order underneath it all.
The bizarre yet tight connections between randomness and order are what all of
Ulam's greatest discoveries are about. His style was iconoclastic, to be sure. He was
perfectly able to do mathematics in the




Mathematical Chaos and Strange Attractors                                            387

classical "theorem-proof-theorem-proof" way, but he delighted in the experimental
approach, using computers to study crazy behaviors of oddball functions he dreamt
up. In some sense, Ulam was a genuine mathematical artist, unlike so many
mathematicians. A piece of math by Ulam often feels much more like a creation than
like a discovery. It is more idiosyncratic, more easily recognizable as the product of a
particular mind, than most mathematical discoveries are.
         Aside from being fascinated by mathematics itself, Ulam was also fascinated
by the human mind's workings, and he strove to express his vague but provocative
intuitions in his writings. I always think of his "ten dogs" theory of memory. The idea
is this: When you are searching for a memory that eludes you but that you know is
there, what you in effect do is release ten "dogs" in your brain and let them go
"sniffing" in parallel. Each dog will start to rummage around here and there,
sometimes going in circles, sometimes smelling down wrong alleys, but since there
are a bunch of them, you can afford to let them smell out many false pathways. They
don't need to be very bright; they just need to have had a whiff of the original idea,
and they will follow that spoor high and low. Eventually, it is likely that one dog or
another will trot home carrying the desired memory in its mouth. Ulam's
autobiography, Adventures of a Mathematician, is packed with such glimmerings
about how minds work, as well as with droll anecdotes about many of this century's
most brilliant mathematicians.
         Ulam was very curious about language. He and his wife came to this country
about 50 years ago, and both loved the English language. But whereas Stan never lost
his strong Polish accent and constantly made errors in English, Francoise eliminated
almost every trace of her French accent and became a virtually flawless speaker,
whose mastery of idiomatic phrases exceeded that of most native speakers. This
caused some amusing light-hearted bickerings between them that I witnessed.
Francoise one day used some baseball idiom such as "he threw them a curve ball" or
"in the ball park", and Stan immediately objected, saying "You can't use that
expression! You didn't grow up playing baseball, so you don't really know what it
means!" Francoise defended herself, saying that she had a good idea of its literal
meaning but that in any case Stan's point was a red herring. I bought her argument
lock, stock, and barrel. After all, how many native speakers of English know what
domains such phrases as "red herring" or "lock, stock, and barrel" come from? Yet we
certainly all use many such phrases and feel perfectly entitled to do so.
         Like many of the brightest mathematicians and physicists working during and
just after World War II, Stan Ulam got involved in military projects. His invention,
with John von Neumann, of the Monte Carlo method was a key element in the
development of the hydrogen bomb. The same forces that drove him to wonder about
the cardinality of abstrusely defined sets and the dimensionality of peculiarly defined
spaces also guided him to accurate ways of modeling the statistics of chain reactions.
At the time he did the work,




Mathematical Chaos and Strange Attractors                                           388

the nature of the dilemma it would lead humanity as a whole into was not so. clear as
it now is. To be sure, Einstein had warned us about our slow drift into unparalleled
peril, but few people had Einstein's clarity of vision. One of the paradoxes about
people is that they are so small compared to the events they can be involved in. Stan
Ulam was an ant in a vast colony, and though his role was more significant than that
of most ants, he still had no control over the nature of the colony itself. Human nature
is one thing, but humanity's nature is another thing.
        A good and generous person like Stan Ulam can still be a part of a bad and
frightful thing like the arms race. Clearly Ulam had many afterthoughts about his role
in these developments, and it is to his credit that he tried to think it all through
rationally. Others in similar positions have been far more trapped and narrow-minded,
unable to see, or to admit seeing, the complex tragedy that has been unfolding as a
consequence of their small actions joined with the small actions of many, many
others. For me it was a privilege to get to know and be friends with this warm and
insightful man. I hope that in the long run, Stan Ulam's contributions to mathematics
will prove to have outweighed his contributions to a potential Armageddon.

                                         *   *   *

        One of the basic themes of this column is what I call locking-in. For no
particular reason, I failed to use that term in the column, but it is a good term. The
imagery I wish to convey is that of a system that seeks and gradually settles into its
own most stable states, and the mechanism whereby it seeks and attains such loci of
stability is feedback. A system that locks into a state is in a stable equilibrium, which
means that if you perturb it somehow, it will swiftly return to the state it was in-there
are restoring forces that push it back. Perhaps the most primordial image is that of the
particle in the potential well-for example, a marble sitting at the bottom of a round
dish. If you ping it lightly with your finger, it will oscillate for a while, but eventually
will come to rest again just where it was before: at the sole stable fixed point of the
system. Here, as in the column, "fixed point" means that the system's "output" at time
t (namely, the marble's position at time t) is identical to the "input" at time t -1
(namely, the marble's position at time t - 1). In this case, the attractor is a single point
in space, so it is ridiculously easy to visualize. Most of the attractors in the chapter,
however, were orbits rather than single points, so they are slightly more abstract.
However, if you think of an orbit as simply a point in a multidimensional space, then
the concept of zeroing in on a fixed point and the concept of settling down in a stable
orbit merge somewhat.
        One of the most intuitive as well as charming examples of locking-in is the
search for a solution to Raphael Robinson's puzzle in Chapter 2:




Mathematical Chaos and Strange Attractors                                              389

   In this sentence, the number of occurrences of 0 is -, of 1 is -, of 2 is \_, of 3 is \_,
   of 4 is -, of 5 is -, of 6 is of 7 is \_, of 8 is \_, and of 9 is -.

One way to search for a solution to this puzzle is to fill in the blanks with an arbitrary
sequence of ten numbers, such as <0,1,2,3,4,5,6,7,8,9>, and see what happens when
you check out the truth of the resulting sentence. It turns out actually to have two
occurrences of each digit. Thus the vector <0,l,2,3,4,5,6,7,8,9> leads to the vector
<2,2,2,2,2,2,2,2,2,2> by the process we'll call "Robinsonizing". Where does that
vector lead? Clearly to < 1,1,11,1,1,1,1,1,1,1 >, which leads to < 1,12,1,1,1,1,1,1,1,1
>, which leads to < 1,11,2,1,1,1,1,1,1,1 >, which leads to < 1,11,2,1,1,1,1,1,1,1 > -and
to and behold, we've entered a closed loop!
        This vector < 1,11,2,1,1,1,1,1,1,1 > is like. a whirlpool or a vacuum cleaner: it
sucks things near to it into its vortex. It is a trap, a fixed point -an attractor. It is not
unique; there is another such vortex, which I will leave it to you to find, Furthermore,
there is at least one two-state loop, or period-two attractor, that I know of. I 'have
reason to suspect that everything leads to one of those three attractors, but I could be
wrong. You could search for a period-two attractor by writing down a vector of length
twenty and generating its successor length-twenty vector as follows: Let the new
vector's first half be derived from the old one's second half by Robinsonizing, and let
the new one's second half be derived from the old one's first half by Robinsonizing. If
you now iterate this double-barreled Robinsonizing operation starting with a random
seed, you will eventually settle down on a fixed point.
        Notice that we are now calling a period-two attractor a "fixed point". Notice
also that this is a "point" in a twenty-dimensional space! The point is, we can view the
system either as bouncing back and forth between two ten-dimensional points (a
period-two attractor) or as sitting still on a fixed twenty-dimensional point. If by
chance there were a loop of length four, we could similarly think of it as being a fixed
point in a 40-dimensional space. As long as we're willing to "up" the dimensionality
of the space, we can store more and more information in a single point. Thus fixed
points and stable orbits are very close concepts.

                                         *   *   *

       This example serves to illustrate how feedback-plugging the system's output
back into the system as input-ushers you to the fixed points. Why should this be so?
Why could the system not thrash about randomly, somehow avoiding all fixed points?
In short, why are fixed points so often attractive? Why could there not be a large
number of fixed points that are totally isolated, like islands in a vast sea, unreachable
via any obvious route? Could there not be fixed-point "anti-whirlpools" that repel any
approacher




Mathematical Chaos and Strange Attractors                                                390

that is not dead on target? In the case of Robinson's puzzle, the answer is no; but there
are such systems. Indeed, in the column I pointed out how there are repellent as well
as attractive fixed points for functions of the form 41x(1-x). But in general, it seems
to be a very good rule of thumb to search for fixed points by starting out somewhere
at random and then hoping that you will get sucked into a stable orbit. Most likely you
will, and you will thereby discover a locus of stability, a locked-in solution.
         Even more remarkable, it seems generally reliable that you are more likely to
be sucked into a short loop than a long one, if short ones exist. Thus, generally
speaking, the stablest behavior of a system seems also to be its simplest behavior.
This is true for systems of nearly any sort one can imagine. In the hydrogen atom, for
instance, the ground state-the lowest-energy state-is spherically symmetric, and is the
only one to have that simple property. Why should this be so, all across the board?
Why are stable things the simplest things as well? Or, conversely, why are the
simplest things the stablest of all? A toughie.

                                        *   *   *

        A puzzle more complex than Robinson's but similar in flavor is the search for
self-documenting or self-inventorying sentences, which was carried out with such
great gusto by Lee Sallows (see Chapter 3). His "logological rocket" was a machine
for seeking attractive fixed points in a certain logological space. The book Loopings
by Aldo Spinelli is a remarkable investigation of regions of a similar logological
space, and his search is guided by the same old principle: that starting somewhere
random and relying on feedback to get you somewhere "better" is the most likely way
to discover a fixed point. This is a most strange way of looking for what might seem
something elusive and precious, yet strange though it might be, it is very robust.
        In Chapter 3's Post Scriptum, I stated that I felt Lee Sallows was overconfident
in wagering that a computer search for a self-documenting sentence beginning "This
computer-generated pangram contains ..." would not succeed in ten years. The reason
is simple. Lee did not consider the idea of "iterative convergence" to a solution-that is,
the idea of Robinsonizing, applied to self-descriptive sentences. You begin with a
sentence of the right form, but where all the numbers are randomly chosen. It's a
blatant lie, but who cares? You just feed it to a program that counts all its letters and
spits out a new sentence with the new letter-counts replacing the old guesses. Around
and around you go ... It is almost certain that you will pretty soon fall into an
attractive orbit. Probably most orbits are fairly lengthy loops, and thus do not yield
self-documenting sentences-but again, who cares? Just try it again with a different
random seed, and keep on doing so, until you find a fixed point.
        This method may sound too simple, but it works. I suggested it to Bob




Mathematical Chaos and Strange Attractors                                             391

French, one of the two translators into French of Godel, Escher, Bach, and he was
gung-ho about implementing such a program. Within a short time, he had one up and
running: He sent me this note about his discoveries:

  I wrote a nice program to solve the Pangram Problem and got an answer,
  written, much to my annoyance, in "francaix". It is:

             Cette phrase contient cinq a, cinq c, troix d, douze e, un f, un g,
     quatre h, treize i, huit n, six o, troix p, six q, huit r, six s, quatorze t, dix u,
     un v, sept x, \& quatre z.

  Unbelievably, in programming it, I had put the wrong goddam spelling of
  "trois" into the program. Oh well, when I corrected the mistake, I didn't get an
  answer immediately, but I'm confident that it'll come, in correctly spelled
  French, when I- get back to work on the thing.

The point is, you don't need to perform a brute-force search through the entire space
of all possible combinations of numbers filling the 26 blanks in order to find a perfect
self-documenting sentence, not by a long shot! A Robinsonizing routine, together
with a simple-minded loop detector, will do the trick quite easily, as long as you're
willing to try a bunch of different seeds. The pulling-power of short loops will
undoubtedly snag you sooner or later, and you'll have found your target sentence!
        My friend Larry Tesler, equally spurred on by Sallows' challenge when it
appeared in print .in A. K. Dewdney's new Scientific American column called
"Computer Recreations" in October 1984, coded up the Robinsonizing method in a
program and soon his computer fell into a loop that seemed very close to a solution.
By changing his program's search technique at that point, Tesler was then easily able
to home in on a winner, which he gleefully sent off to both Dewdney and Sallows.
Tesler's sentence runs as follows:

  This computer-generated pangram contains six a's, one b, three c's, three d's,
  thirty-seven e's, six f's, three g's, nine h's, twelve i's, one j, one k, two l's, three
  m's, twenty-two n's, thirteen o's, three p's, one q, fourteen r's, twenty-nine s's,
  twenty-four t's, five u's, six Vs, seven w's, four x's, five y's, and one z.

                                         *   *    *

       Locking-in is perfectly illustrated by the hypothetical book Reviews of This
Book, described in Chapter 3. There I characterized the method of its creation as
resembling the construction of "self-consistent" solutions via the "Hartree-Fock"
method. What does that mean? It boils down to the same thing once more. It turns out
to be very hard-in fact, impossibleto give closed-form solutions to the equations
describing any atom more complicated than a hydrogen atom, with its single electron.
When you have three bodies, as in the helium atom with its two electrons and a
nucleus, the mathematical complexity is overwhelming. The problem is in essence
that




Mathematical Chaos and Strange Attractors                                                   392

each electron would "like" to be in a simple hydrogen-like state around the nucleus,
but the other one is blocking it from so doing. How can they "cooperate" with each
other to find a stable mode of coexistence?
        One way to study this mathematically, suggested first in 1928 by the English
physicist Douglas Rayner Hartree, is to try to converge on a good description of the
total system by starting out with a false solution-a mathematical description of a state
known to be wrong, but easy to describe. (For instance, you could pretend that both
electrons are in simple hydrogen-like states.) Then you see how each electron
"perturbs" the other one out of the presumed state it was in. This leads you to a
different-and probably no less fictitious-state. But at least you've made progress, in
that you've taken into account the "first-order" effects each electron would have on
the other one. Now you do the same thing over again-that is, you see how the
perturbed states would perturb each other. This gives you "second-order" corrections-
and so on and so on. Eventually-and this is the beauty of the method-the starting point
of your calculations gets totally buried, and the state converges to what is called a
"self-consistent" solution, very much like the solutions to Robinson's puzzle. What I
mean by saying the starting point gets "buried" is that no matter where you start,
you'll wind up at the same eventual solution-a fixed point, where further iteration has
no effect. In this solution, the two electrons are in equilibrium with each other and do
not perturb each other. And presto-one has "solved" the helium atoml
        Of course, this type of solution is numerical, not analytic: there are no exact
formulas that come out, only numbers. Nonetheless, that's good enough for most
practical purposes. The Russian physicist Vladimir Fock later made a suggestion for
improving the validity of this method of calculation, which involves taking into
account the fact that electrons obey the Pauli exclusion principle, a complication that
Hartree had ignored. That is the reason for the hyphenated name; however, Hartree is
the inventor of the general principle of calculating self-consistent solutions for many-
body systems.

                                        *   *   *

        This idea of locking-in recurs throughout science. In Gödel, Escher, Bach, I
discussed the phenomenon called renormalization-the way that elementary particles
such as electrons and positrons and photons all take each other into account in their
very core. The notion is a mathematical one, but for a good metaphor, recall how your
own identity depends on the identities of your close friends and relatives, and how
theirs in turn depends on yours and on their close friends' and relatives' identities, and
so on, and so on. This was the image I described for "I at the Center" in the Post
Scriptum to Chapter 10. Another good graphic representation of this idea is shown in
Figure 24-4, where identity emerges out of a renormalization process.
        The tangledness of one's own self is a perfect metaphor for




Mathematical Chaos and Strange Attractors                                             393

understanding what renormalization is all about. And the best way to imagine how
you emerge from such a complex tangle is to begin by imagining yourself as a
"zeroth-order person"-that is, someone totally unaware and inconsiderate of all others.
(Of course, such a person would be barely a person, barely a self at all: a perfect
baby.) Then imagine how "you" would be modified if you started to take other people
into account, always considering others as perfect babies, or zeroth-order people. This
gives a "first-order" version of you. You are beginning to have an identity, emerging
from this modeling of others inside yourself. Now iterate: second-order people are
those who take into account the identities of first-order people. And on it goes. The
final result is renormalized people: people who take into account the identities of
renormalized people. I know it sounds circular, and indeed it is, but paradoxical it is
not-at least no more than are the fixed points of Raphael Robinson's puzzle!
"Circular" is not synonymous with "paradoxical", although many people mistakenly
assume it is. We shall re-encounter this notion of renormalized people in Chapter 30
and beyond, where it will in fact clear up some seeming paradoxes involving
cooperation and egoism.
         This close connection of locking-in to the deepest essence of personhood plays
a central role also in Chapters 22 and 25, where "who" one is is portrayed as emerging
from a "level-crossing feedback loop", in which a sophisticated perceiving system
perceives limited aspects of its own nature, and by feeding them back into the system
creates a type of locking-in. The locked-in loop itself is given a name, and that name,
for every such system, is "I".
         The idea of a system with an I, watching its own behavior, is closely related to
the wellsprings of creativity (recall the cycle underlying creativity discussed in the
Post Scriptum to Chapter 12, and that to Chapter 10 as well). We will delve into this
in depth again in Chapter 23, trying to come to grips with another seeming paradox:
that of mechanizing what seems by definition to be nonmechanical and
nonmechanizable-the creative act. Once again we'll see vicious paradox dissolve into
benign cycles.
         In short, locking-in-that is, convergent and self-stabilizing behaviorwill surely
pervade the ultimate explanation of most mysteries of the mind. One example is the
question of memory retrieval. How do things that are only vaguely similar to each
other stir up rumblings of recollection, and eventually trigger the retrieval of
amazingly deep abstract resemblances? One theory, best formulated and articulated
by cognitive scientist Pentti Kanerva of Stanford University, sees the initial input as a
seed-a vector in a very high-dimensional space, analogous to the seed vector that we
fed into the Robinsonizing machine. The seed is fed into memory-retrieval
mechanisms, which convert it into an output vector that is then fed back in again. This
cyclic process continues until it either converges on a stable fixed point-the desired
memory trace-or is seen. to be wandering erratically without any likelihood of locking
in, tracing out a chaotic




Mathematical Chaos and Strange Attractors                                             394

sequence of "points" in mind-space. The details of how this is accomplished in
Kanerva's beautiful theory are beyond the scope of this book, but this "self-
propagating search" provides another remarkable example of the many ways that
locking-in can be exploited.
        Closely related to memory retrieval is the problem of perception, or pattern
recognition. As I mentioned in the Post Scriptum to Chapter 4, this central aspect of
mind has been best modeled on computers in programs whose strategy is similar to
that of Kanerva's model: there is a superficial sweep that narrows the field somewhat,
followed by a deeper sweep that narrows it further, and so on (the "terraced scan" I
described in the postscript to Chapter 5). This bottom-up processing is complemented
by concurrent top-down processing driven not by the input, but by expectations of
what is "out there" to be recognized. The swirling activity in which bottom-up and
top-down processes seek a reconciliation with each other leads to a gradual kind of
"crystallization", in which many small pieces of evidence align with, and mutually
reinforce, each other. The ultimate justification for some of them resides, of course, in
the raw perceptual input, while for others of them it resides in the richness of previous
experiences stored in memory. The combination of all these mutually confirming
hypotheses results in a globally optimal interpretation of the input: an act of
recognition. Once again, locking-in carries the day.
        One final example of locking-in is the subject of Chapter 27: the question of
the inevitability (or evitability) of the genetic code. This central question about the
molecular foundations of life turns out to revolve about two distinct senses of the
word "arbitrary". I shall let that Chapter speak for itself, however.

                                       *   *   *

        In the Introduction, I described the space of my columns as gradually
emerging as, month by month, I revealed one more dot in that space. What is this, if
not a Poincare map of my mental meanderings? During my column-writing era, my
mind would light up like a monthly firefly and reveal where it was to the outside
world! I just wonder: Would the shape I was thus tracing out turn out to be a strange
attractor?
        It seems appropriate that at this midpoint of the book, we have identified a
unifying theme-or rather, thema, to be more faithful to the title. Locking-in seems to
be a key to the metamagics of Snarls, of Society, of Slipping ... of Strangeness, of
Substrate, of Stability ... of Survival.




Mathematical Chaos and Strange Attractors                                            395

